\documentclass[10pt]{article}
\usepackage[margin=1.5cm]{geometry}
\usepackage{graphicx,booktabs,amsmath,amssymb}
\usepackage[hidelinks]{hyperref}
\usepackage{fancyvrb,amsthm}
\usepackage{xr}
\externaldocument[M-]{../manuscript/main}
\newtheorem{propS}{Proposition}
\renewcommand{\thepropS}{S\arabic{propS}}
\newcommand{\PP}{\mathbb{P}}
% R20 compact review candidate generated by make_compact_supplement.py.
% Proofs/table inputs are frozen R19 copies; tables below use frozen saved JSON.
\renewcommand{\thetable}{S\arabic{table}}
\renewcommand{\thefigure}{S\arabic{figure}}
\renewcommand{\theequation}{S\arabic{equation}}
\title{Supplementary Material for ``Innovation-Normalized Detection in Compound-Gaussian Clutter: Exact Laws, Conformal Thresholds and Certified Integration under Pulsed Interference''}
\author{M. U. Raza, S. Ahmed, A. Ahmed, M. Atif Shahzad, and Syed M. Kazam Abbas Kazmi}
\date{}

\begin{document}
\sloppy\maketitle
\noindent This supplement provides proofs of the paper's conditional laws and supporting results on calibration, target visibility, interference and real targets. Numerical tables come from frozen saved outcomes; the registered extensions below add new outcomes to the preserved studies. Protocols, saved summary outputs and plotting arrays are available in the accompanying repository (release \texttt{r22-scope-submission}, \texttt{research/r7c}); Section~\ref{sup:repository} gives their paths. Study labels identify internal hash-recorded protocols, not independent experiments or an external registry. Failed conditions and post-hoc choices remain explicit.

\section{Data exposure}
Table~\ref{tab:exposure} states which data each stage used, so that the strength of each validation claim can be judged. Only the hold-out series, the 77~GHz data, NetRAD and the new UW pedestrian records were not used while the methods were being developed. The long-trajectory and migration extensions reuse exposed recordings.
\begin{table}[h]
\centering\small
\caption{Data exposure by stage. ``Seen'' means that outcomes on these data were computed while methods, orders or comparisons could still change.}
\label{tab:exposure}
\begin{tabular}{p{4.3cm}p{5.2cm}p{6.2cm}}
\toprule
Data & Used for & Status of the reported results\\
\midrule
IPIX, 14 sessions, rotation 1 & Method development (R1--R7): method list, model orders, comparisons & Seen; not reported\\
IPIX, same sessions, rotations 2--3 & Main study, detection, onset, baselines, classical detectors, order-statistic scores (R7--R11) & Frozen-protocol evaluation, computed once per protocol; the recordings were seen in rotation 1, so this is not an untouched confirmation\\
IPIX, day transfer & Transfer and ACI analyses & Frozen-protocol evaluation on the same recordings\\
IPIX \#269, \#287 & Hold-out & Hashed and unopened until the main analysis was complete; evaluated once (small test sets)\\
JKU 77~GHz, reference run & Second-radar study & Not used in development; frozen protocol; the general-form gain curve is post hoc\\
JKU 77~GHz, interference runs A and C & Real-interference test (R11); certified integration (R12) & R11: only chirp-hit statistics inspected before its protocol. R12: R11 outcomes on these runs had been seen\\
IPIX, rotations 2--3 & Guarded scale, clip level, trimming, blanking (R14) & Frozen-protocol evaluation, computed once; same recordings as R7--R12\\
IPIX, primary target cells & Real-target study (R12) & Excluded from the R7--R11 clutter studies; a development check (V4, before R7) had run IN-ARCP on the primary cell and found it blind. Frozen protocol, evaluated once\\
NetRAD S-band, node 3, 14 recordings & Second-radar sea-clutter confirmation (R15) & Untouched: protocol frozen before any statistic of the samples was computed; evaluated once\\
JKU 77~GHz, pedestrian run (scenario A) & Real onsets of a walking person (R17) & Untouched: protocol frozen and pushed before the file was downloaded; evaluated once; descriptive (E1 not met)\\
UW 77~GHz, three moving-person records & Detector-external camera association & Metadata-selected before new FFTs or detector outputs; frozen protocol; all correlation gates fail and the primary clean-entry subset is empty\\
\bottomrule
\end{tabular}
\end{table}

\section{Protocol deviations}\label{sup:deviations}
\begin{itemize}
\item Main IPIX study: before any run, an encoding error in the protocol file was repaired byte-exactly and the hash rewritten; no outcome had been computed. A GPU backend was added for the likelihood evaluation of the noise-aware fit; it gives identical fitted parameters to the CPU path.
\item Second radar: the descriptive hit rule summed range bins 20--247 instead of 20--250, because the fixed preprocessing removes bins above 247. The general-form gain curve was chosen after the pre-specified matched form failed to capture the magnitude of the measured gain, so it is post hoc.
\item Order-statistic study (R11): the protocol's rationale counted $j+1$ target-corrupted history innovations at look $j$ after onset; the correct count, $j$, was adopted after the data were seen. The data follow the corrected count, and the main text uses only the corrected statement.
\item Final study (R12): before the real-data runs, the numerical search for the Fisher's $g$ threshold was re-bracketed; the law is unchanged.
\end{itemize}

\section{Unmet expectations}\label{sup:unmet}
Each expected value was written down before the outcome was computed. For the earlier R12, R14, R15 and R17 protocols, thirteen of 53 expected outcomes were not met and four were left without a verdict. The long-trajectory, migration and camera studies report their outcomes in their own sections.
\begin{itemize}
\item Pedestrian study (R17; E1 not met, E2--E5 without a verdict): the onsets were expected in at least eight distinct five-frame blocks and fell in seven, because the person was tracked for only about six seconds. Every other result is therefore descriptive; all lie on the side the protocol expected (Section~\ref{sup:r17}).
\item Final IPIX study (R12; 12 of 18 met): calibration of the saturating clipped statistic; the Tyler law (expected within twofold); the CA power detector's white-noise law (expected off more than threefold); the cost of clipping (expected at most 1.5~dB); false alarms of non-coherent integration under real FMCW interference (expected at least fivefold in both runs); and the gain of AR gap reconstruction (expected at least 1~dB).
\item Improvement study (R14; 12 of 15 met): the guard's onset cost, the certificate at $10^{-3}$ and the cost of trimming (Section~\ref{sup:r14}).
\item NetRAD confirmation (R15; 12 of 15 met): calibration of six statistics at $10^{-4}$, the pivotality ranking against CA16, and a rise of false alarms in segments with flagged pulses (Section~\ref{sup:r15}).
\item Main IPIX study ($m=16$, $\alpha=0.1$; radius relative to IN-ARCP): the noise-aware gain was expected to concentrate in noise-affected sessions. It did for NA-AR(1) (0.915 noisy against 0.989 clean) but not for NA-AR(2) and NA-AR(4) (0.871 and 0.866 noisy against 0.739 and 0.709 clean), whose gain comes mainly from the AR clutter dynamics (expectation Q1, mixed).
\item Onset study: the gain of IN-ARCP over the long-window power detector was expected to shrink below 3~dB by an 8-pulse dwell. It persisted: 8.4~dB [5.2, 11.1] at $P_{\rm d}=0.5$ (expectation O3, refuted).
\item 77~GHz radar: the noise-aware scores were expected to be flatter than IN-ARCP across CNR classes. The AR(4) version was (mean largest class deviation 0.055 against 0.066); the AR(1) version was not (0.067; expectation J2, mixed).
\item Classical detectors (R10). The frozen \texttt{r10/r10\_summary.txt} prints expectation R1 as FAILED. It is met in substance. For a target that rises over $L$ pulses, reaching full amplitude at the first look, IN-ARCP's first-look gain over the long-window power detector (CAloc in the paper) falls from 9.86 to 8.65 to 4.65~dB for $L=1,4,8$. At $L=16$, CAloc reaches $P_{\rm d}=0.5$ at 9.1~dB while IN-ARCP stays below it up to 25~dB, so the gain is below $-15.9$~dB, well under the 5~dB bound. PAMF-H was expected to need at least 3~dB less SCR than integrated IN-ARCP over an 8-pulse dwell at $P_{\rm d}=0.5$; it needed 1.0~dB [0.3, 1.7] less (4.0~dB at $P_{\rm d}=0.9$). Gated ACI was expected to keep the clean false-alarm rate within $[0.8,1.2]\alpha$ at 5\% target contamination; it gave $0.55\alpha$ ($\gamma=0.02$) and $0.68\alpha$ ($\gamma=0.005$).
\item Order-statistic study (R11; 2 of 11 met): quintile false-alarm spread of the median OS score at most 2.0 (2.06); its cost for abrupt targets at most 1.5~dB (1.99~dB [0.89, 3.39]; Corollary~3 gives 0.69~dB); a first-look gain of at least 3~dB over IN-ARCP for 8-pulse ramps (1.46~dB); the integrated OS score closing half of integrated IN-ARCP's gap to PAMF-H at $P_{\rm d}=0.9$ (it was 2.5~dB worse); OS-law coverage error at most 0.015 (0.021); on the dense real-interference run, the OS score halving IN-ARCP's masking (drops 0.39 and 0.35), zeroing cutting the hit-pulse false-alarm rate by half (39\%), and the OS score after zeroing within 0.5~dB of IN-ARCP (0.71~dB worse); on the reference run, IN-ARCP beating CA16 (CA16 was 0.78~dB better, the loss Proposition~2 predicts for $c<1$).
\end{itemize}


\section{Proof of Proposition 2 and the latent-texture Mondrian result}
Equation, proposition and theorem numbers without the prefix S refer to the paper.

\emph{Proposition~2.} If $e\sim\mathcal{CN}(0,v)$ and a Swerling-1 target of power $SP$ is added, a threshold test on $|e|^2/v$ at false-alarm probability $\alpha$ detects with probability $\alpha^{1/(1+SP/v)}$, while the power-only test on $|Y|^2/P$ detects with $\alpha^{1/(1+S)}$; equal $P_{\rm d}$ therefore requires SCRs in the ratio $G=P/v$. For the center $rH_m$ under \eqref{M-eq:r7model} with AR(1) speckle, $v=\nu\{c(1+|r|^2-2\Re\,\bar r\rho)+1+|r|^2\}$, which gives $G_{\rm IN}$. For the optimal predictor from the infinite past, $v$ is the one-step prediction-error variance of the AR(1)-plus-white-noise process, $\exp\{\int\log f\}$ by the Kolmogorov--Szeg\H{o} formula \cite{S-brockwell1991}. With $f(\omega)=\nu\{c(1-|\rho|^2)+|1-\rho e^{-j\omega}|^2\}/|1-\rho e^{-j\omega}|^2$, the denominator contributes zero, and the numerator factorizes as $\nu V|1-\zeta e^{-j\omega}|^2$ with $V(1+|\zeta|^2)=\xi$, $V\zeta=\rho$ and $|\zeta|<1$, which gives $v=\nu V$ and $G_{\rm opt}$. Because the optimal predictor's error variance is at most that of any linear predictor, including $0$ and $rH_m$, $G_{\rm opt}\ge\max\{1,G_{\rm IN}\}$.

\emph{Latent-texture Mondrian calibration.} Mondrian (group-conditional) conformal prediction calibrates separately within bins of an observed statistic and is valid conditional on the observed bin; Dewolf \emph{et al.}\ report that estimated taxonomies mix up the true classes \cite{S-dewolf2025}. For a latent texture the consequence can be stated exactly.
\begin{propS}\label{prop:mondrian}
Let $\nu=0$ and $X=\sigma Z$, with $\sigma>0$ distributed as $G$, non-degenerate and independent of $Z\sim\mathcal{CN}(0,R)$. Let $S(Z)$ be a scale-invariant score with a continuous distribution function $F_S$ that is strictly increasing on its support, let $T(H)>0$ be continuous and homogeneous of degree one, and calibrate within fixed bins $B_b=[e_{b-1},e_b)$, $0=e_0<\dots<e_K=\infty$, $K\ge2$, of $T(H)$ with population quantiles $q_b$. Let $\mathcal C(\sigma)$ be the coverage conditional on the texture.
(i) If $S(Z)$ is independent of $T(Z_H)$, then $\mathcal C(\sigma)\equiv1-\alpha$.
(ii) If $S(Z)$ given $T(Z_H)=t$ is strictly stochastically decreasing in $t$, then $\lim_{\sigma\to\infty}\mathcal C(\sigma)=F_S(q_K)<1-\alpha<F_S(q_1)=\lim_{\sigma\to0}\mathcal C(\sigma)$, whereas the unbinned threshold gives a constant $\mathcal C(\sigma)$. Condition (ii) holds exactly for IN-ARCP at the true coefficient with $T=s_\rho$.
\end{propS}
\emph{Proof.}
Write $t=T(Z_H)$, so that the observed statistic is $T(X_H)=\sigma t$ and $S(X)=S(Z)$. For bin $b$ let $w_b(t)=\mathbb{P}\{\sigma t\in B_b\}$. Since $\sigma$ is independent of $Z$, for every $x$
\begin{equation}
\mathbb{P}\{S\le x,\ \sigma t\in B_b\}=\mathbb E\big[\varphi_x(t)\,w_b(t)\big],
\tag{S1}
\end{equation}
where $\varphi_x(t)=\mathbb{P}\{S\le x\mid t\}$.
(i) If $S$ is independent of $t$, then $\varphi_x\equiv F_S(x)$, every bin-conditional law of $S$ equals $F_S$, all $q_b$ equal $F_S^{-1}(1-\alpha)$ and $\mathcal C(\sigma)=1-\alpha$.

(ii) Let $x_0=F_S^{-1}(1-\alpha)$. For the top bin, $w_K(t)=\mathbb{P}\{\sigma\ge e_{K-1}/t\}$ is nondecreasing in $t$. Because $T$ is continuous, positive and homogeneous and $Z_H$ has full support, $t$ has full support on $(0,\infty)$; as $e_{K-1}/t$ then ranges over $(0,\infty)$ and $G$ is non-degenerate, $w_K(t)$ is not almost surely constant. By assumption $\varphi_{x_0}$ is strictly increasing. Let $t'$ be an independent copy of $t$, $D_\varphi=\varphi_{x_0}(t)-\varphi_{x_0}(t')$ and $\Delta_w=w_K(t)-w_K(t')$. Then
\[
\operatorname{Cov}\{\varphi_{x_0}(t),w_K(t)\}=\tfrac12\,\mathbb E[D_\varphi\Delta_w]>0,
\]
because $D_\varphi\Delta_w\ge0$, with strict inequality on the event $\{w_K(t)\ne w_K(t')\}$, which has positive probability. By (S1),
$\mathbb{P}\{S\le x_0\mid\sigma t\in B_K\}=\mathbb E[\varphi_{x_0}w_K]/\mathbb E[w_K]>\mathbb E[\varphi_{x_0}]=1-\alpha$.
The bin-conditional distribution function of $S$ is continuous, so there is $x_1<x_0$ at which it still exceeds $1-\alpha$; hence $q_K\le x_1<x_0$ and, $F_S$ being strictly increasing, $F_S(q_K)<1-\alpha$. For the bottom bin, $w_1(t)=\mathbb{P}\{\sigma<e_1/t\}$ is nonincreasing, the covariance is negative, and the same steps give $q_1>x_0$ and $F_S(q_1)>1-\alpha$.

The coverage at texture $\sigma$ is $\mathcal C(\sigma)=\sum_b\mathbb{P}\{S\le q_b,\ \sigma t\in B_b\}$. As $\sigma\to\infty$, $\sigma t\to\infty$ almost surely, so $\mathbf 1\{\sigma t\in B_K\}\to1$ and, by dominated convergence, $\mathcal C(\sigma)\to F_S(q_K)$; as $\sigma\to0$, $\mathcal C(\sigma)\to F_S(q_1)$. An unbinned threshold $q$ gives $\mathcal C(\sigma)=F_S(q)$ for every $\sigma$.

Finally, for IN-ARCP at $r=\rho$ and $T=s_\rho$, $S=|E|/t$ with the innovation $E$ independent of $Z_H$, so $\varphi_x(t)=\mathbb{P}\{|E|\le xt\}$ is strictly increasing in $t$, which is condition (ii). \hfill$\blacksquare$

\section{Zero-noise limit of the noise-aware score}
\begin{propS}\label{prop:reduce}
Let $R$ be AR(1) with coefficient $\rho$ and fix a nonzero history $h$. As $\nu\downarrow0$, $\nu\widehat c(h)\to h^{\mathsf H}R_H^{-1}h/m$, $\mu(h)\to\rho h_m$ and $s(h)\to s_\rho(h)$, so the NA score converges to the IN-ARCP score at $r=\rho$.
\end{propS}
\emph{Proof sketch.}
With $\gamma=\nu c$ the history log-likelihood converges locally uniformly on $\gamma>0$ to $-m\log\gamma-\gamma^{-1}h^{\mathsf H}R_H^{-1}h$, whose unique maximizer is $h^{\mathsf H}R_H^{-1}h/m$; coercivity at $0$ and $\infty$ gives convergence of the maximizers. The Markov property gives $\mu\to\rho h_m$, and $(1-|\rho|^2)R_H^{-1}=W_\rho^{\mathsf H}W_\rho$, with $W_\rho$ the bidiagonal whitening matrix, identifies the limiting predictive variance with $s_\rho(h)^2$.


\input{sec_r12_proofs.tex}
\section{Numerical verification and frozen-output corrections}\label{sup:indep}\label{sup:r12}
The saved \texttt{r12/verify\_laws\_output.txt} contains 400,000-replication Monte Carlo checks. Its 65 primary checks are the rows with $z$-scores excluding bracketed labels: two compound-Gaussian tests of the Gaussian Kraut--Scharf law and four asymptotic fixed-point-law checks. Every primary check lies within three standard errors (maximum $|z|=2.9125$, printed as 2.91); the fixed-point checks have $|z|<1.6$. Beta and certificate rows report distribution agreement or violation rates, not $z$-scores. The compound-Gaussian Kraut--Scharf tests fail 2.8--3.4-fold at $10^{-2}$; this failure remains part of the evidence.

The independent \texttt{r14/theory\_check\_independent.py} imports no implementation of the manuscript laws. Its saved output checks repeated eigenvalues, singular scales, the post-onset law, horizon, tested strong-target OS regime, integration, known-bin P-ANMF, Beta calibration, and certificates under strong, matched and cancelling amplitudes. Bursts calibrated as independent hits violate design (clipped $P_{\rm fa}=0.0114$, binary 0.0192 at design 0.01). Matched burst certificates record no exceedance and lose sensitivity. The separate guarded-AR check has maximum $|z|=2.56$ with 400,000 episodes per point. Protected looks use unchanged tested energy and zero target energy in the history; only the contaminated interval shifts by the guard.

\emph{Two corrections accompany the unchanged full logs.} V8 calibration sizes in \texttt{verify\_laws\_output.txt} came from a faulty search. \texttt{r12/calibration\_size.txt} supersedes them: first acceptable sizes 149, 1,497 and 14,978; permanent cutoffs 313, 3,146 and 31,477 at $\alpha=10^{-2},10^{-3},10^{-4}$. The proof supplies the tail completion. The independent log's $n=23,000$ line checks a formula, not the actual set size. Actual per-look calibration sets contain 17,164--23,130 episodes; their i.i.d. Beta-tail probabilities are about 0.02--0.09. Non-overlap does not imply independence. The R10 ramp log prints R1 as FAILED because a numerical crossing is unavailable when IN-ARCP never reaches $P_{\rm d}=0.5$ by 25~dB; the recorded curves meet that expectation in substance, as Section~\ref{sup:unmet} explains.

\subsection{Closed forms against Monte Carlo: all six panels}\label{sup:theoryfig}
Figure~\ref{fig:sup_theory} checks the post-onset, horizon, order-statistic, integration, Doppler and certificate calculations.
\begin{figure}[htbp]
\centering\includegraphics[width=7.16in]{fig_theory_full.pdf}
\caption{Closed forms (lines) against Monte Carlo (markers; $10^5$ episodes per point in (a), (c) and (f), $5\times10^4$ in (b) and (d), $10^4$ to $2.5\times10^4$ in (e)) at the true coefficient, AR(1) compound-Gaussian clutter with $\rho=0.93e^{-0.2\mathrm j}$, lognormal texture, $m=16$, $P_{\rm fa}=0.01$. Panels (a)--(c) and (f) are those of Fig.~\ref{M-fig:theory} of the paper, where (f) appears as its panel (d) and (c) here adds the RMS-scale law for reference. The two panels shown only here are (d), the integration law \eqref{M-eq:dwellpd}, and (e), P-ANMF with unknown Doppler. Doppler offsets are relative to clutter. In (f), interference is 30~dB above local power; curves join Monte Carlo estimates. Open downward triangles mark no false alarm in $10^5$ episodes and are drawn at the resolution $10^{-5}$.}
\label{fig:sup_theory}
\end{figure}


\section{False-alarm calibration and pivotality}
Tables~\ref{tab:ipix_ci} and~\ref{tab:netrad_ci} preserve every Part-L ratio and cluster interval, including the residual bootstrap. They are measured pooled rates, not guarantees of exchangeability or equal rates in individual cells. IPIX intervals resample six days; NetRAD intervals resample fourteen recordings. These intervals are approximate. Saturation and ties can make clipped thresholds conservative.


% Generated from r12/r12_summary.json, Part L. All 13 statistics and residual-bootstrap rows retained.
\begin{table}[htbp]
\centering\small
\caption{IPIX pooled $P_{\rm fa}/\alpha$ [95\% day-cluster interval]. Bold point estimates lie outside $[0.5,2]$. Model rules are those named in the paper's low-rate table.}\label{tab:ipix_ci}
\begin{tabular}{lcc}
\toprule
Statistic & Conformal & Model-based rule\\
\midrule\multicolumn{3}{c}{$\alpha=10^{-2}$}\\
ANMF-Tyler & 1.04 [0.98, 1.07] & 1.62 [1.39, 1.82] \\
ANMF-SCM & 1.06 [0.99, 1.12] & \textbf{2.02} [1.69, 2.28] \\
Power CA & 1.00 [0.99, 1.01] & 0.80 [0.71, 0.85] \\
Clipped, OS scale & 1.05 [0.99, 1.08] & \textbf{4.84} [3.72, 6.72] \\
Integration, AR(4) & 1.02 [0.98, 1.07] & \textbf{2.02} [1.73, 2.27] \\
IN-ARCP & 1.01 [0.99, 1.02] & 0.96 [0.92, 1.02] \\
IN-AR(4) & 1.00 [0.99, 1.01] & 1.03 [1.01, 1.05] \\
NA-AR(4) & 1.00 [0.95, 1.08] & 1.56 [1.40, 1.70] \\
OS scale, $k=8$ & 1.04 [1.01, 1.08] & \textbf{2.15} [1.06, 3.57] \\
Power OS & 1.01 [0.96, 1.04] & \textbf{2.07} [1.47, 2.55] \\
P-ANMF, maximum & 1.04 [0.95, 1.11] & \textbf{4.36} [3.58, 5.29] \\
P-ANMF, one bin & 1.10 [1.01, 1.19] & 1.78 [0.73, 2.61] \\
PAMF-H, one Doppler & 0.99 [0.96, 1.03] & 1.97 [1.06, 2.70] \\
IN-ARCP, residual bootstrap & --- & 1.63 [1.47, 1.75] \\
\midrule\multicolumn{3}{c}{$\alpha=10^{-3}$}\\
ANMF-Tyler & 1.05 [0.96, 1.12] & \textbf{2.66} [1.76, 3.48] \\
ANMF-SCM & 1.04 [0.94, 1.14] & \textbf{3.57} [2.38, 4.54] \\
Power CA & 1.00 [0.94, 1.06] & 0.56 [0.41, 0.63] \\
Clipped, OS scale & 0.67 [0.34, 0.95] & \textbf{17.19} [11.68, 25.51] \\
Integration, AR(4) & 1.05 [0.86, 1.20] & \textbf{4.35} [3.16, 5.33] \\
IN-ARCP & 1.03 [1.01, 1.04] & 1.02 [0.91, 1.12] \\
IN-AR(4) & 1.02 [0.96, 1.08] & 1.19 [1.01, 1.31] \\
NA-AR(4) & 1.05 [1.00, 1.08] & \textbf{2.79} [2.38, 3.24] \\
OS scale, $k=8$ & 1.02 [0.89, 1.13] & \textbf{8.34} [2.00, 15.87] \\
Power OS & 0.98 [0.91, 1.02] & \textbf{5.56} [3.15, 7.36] \\
P-ANMF, maximum & 1.04 [0.90, 1.16] & \textbf{13.06} [10.39, 17.19] \\
P-ANMF, one bin & 1.14 [0.95, 1.26] & \textbf{3.05} [1.01, 5.06] \\
PAMF-H, one Doppler & 1.00 [0.89, 1.08] & \textbf{3.66} [1.66, 5.60] \\
IN-ARCP, residual bootstrap & --- & \textbf{2.43} [2.15, 2.67] \\
\midrule\multicolumn{3}{c}{$\alpha=10^{-4}$}\\
ANMF-Tyler & 0.82 [0.67, 0.95] & \textbf{6.01} [2.87, 9.15] \\
ANMF-SCM & 0.85 [0.75, 1.01] & \textbf{8.07} [3.98, 12.11] \\
Power CA & 0.75 [0.68, 0.90] & \textbf{0.49} [0.26, 0.70] \\
Clipped, OS scale & \textbf{0.03} [0.00, 0.11] & \textbf{59.89} [40.97, 85.35] \\
Integration, AR(4) & 1.02 [0.59, 1.31] & \textbf{14.25} [7.84, 19.14] \\
IN-ARCP & 0.73 [0.54, 0.99] & 1.92 [1.26, 2.61] \\
IN-AR(4) & 0.86 [0.64, 1.11] & \textbf{2.73} [1.48, 3.67] \\
NA-AR(4) & 1.04 [0.84, 1.19] & \textbf{7.46} [4.55, 9.73] \\
OS scale, $k=8$ & 0.92 [0.76, 1.14] & \textbf{45.91} [7.96, 90.46] \\
Power OS & 0.75 [0.57, 1.07] & \textbf{18.33} [8.92, 25.59] \\
P-ANMF, maximum & 1.47 [0.75, 2.57] & \textbf{49.61} [35.38, 73.34] \\
P-ANMF, one bin & 1.10 [0.79, 1.31] & \textbf{5.47} [1.35, 9.99] \\
PAMF-H, one Doppler & 0.89 [0.55, 1.22] & \textbf{9.88} [4.13, 15.83] \\
IN-ARCP, residual bootstrap & --- & \textbf{4.27} [3.44, 4.74] \\
\bottomrule
\end{tabular}
\end{table}


% Generated from r15/r15_summary.json, Part L. All 13 statistics and residual-bootstrap rows retained.
\begin{table}[htbp]
\centering\small
\caption{NetRAD pooled $P_{\rm fa}/\alpha$ [95\% recording-cluster interval]. Bold point estimates lie outside $[0.5,2]$. Model rules are those named in the paper's low-rate table.}\label{tab:netrad_ci}
\begin{tabular}{lcc}
\toprule
Statistic & Conformal & Model-based rule\\
\midrule\multicolumn{3}{c}{$\alpha=10^{-2}$}\\
ANMF-Tyler & 1.01 [0.98, 1.03] & \textbf{2.01} [1.97, 2.04] \\
ANMF-SCM & 1.00 [0.98, 1.02] & \textbf{2.27} [2.20, 2.33] \\
Power CA & 1.01 [0.99, 1.03] & \textbf{0.32} [0.27, 0.36] \\
Clipped, OS scale & 1.05 [1.02, 1.10] & \textbf{2.04} [1.89, 2.23] \\
Integration, AR(4) & 1.06 [1.00, 1.15] & 1.44 [1.36, 1.56] \\
IN-ARCP & 1.04 [0.99, 1.09] & 0.80 [0.63, 0.94] \\
IN-AR(4) & 1.03 [1.00, 1.06] & 0.98 [0.95, 1.02] \\
NA-AR(4) & 1.08 [1.00, 1.19] & 1.90 [1.75, 2.03] \\
OS scale, $k=8$ & 1.05 [0.99, 1.14] & 0.91 [0.73, 1.07] \\
Power OS & 1.01 [0.99, 1.02] & 0.65 [0.60, 0.70] \\
P-ANMF, maximum & 1.04 [0.97, 1.12] & \textbf{2.46} [2.29, 2.66] \\
P-ANMF, one bin & 0.98 [0.90, 1.09] & \textbf{2.28} [1.89, 2.66] \\
PAMF-H, one Doppler & 1.03 [0.97, 1.10] & \textbf{2.38} [2.05, 2.72] \\
IN-ARCP, residual bootstrap & --- & 0.82 [0.55, 1.05] \\
\midrule\multicolumn{3}{c}{$\alpha=10^{-3}$}\\
ANMF-Tyler & 1.01 [0.94, 1.12] & \textbf{3.00} [2.82, 3.21] \\
ANMF-SCM & 1.02 [0.95, 1.11] & \textbf{3.52} [3.29, 3.75] \\
Power CA & 1.07 [1.03, 1.12] & \textbf{0.18} [0.12, 0.23] \\
Clipped, OS scale & 1.13 [1.05, 1.22] & \textbf{4.40} [3.63, 5.45] \\
Integration, AR(4) & 1.36 [0.96, 1.99] & \textbf{2.42} [1.88, 3.22] \\
IN-ARCP & 1.12 [0.99, 1.26] & 0.90 [0.63, 1.13] \\
IN-AR(4) & 1.08 [0.97, 1.20] & 1.06 [0.95, 1.24] \\
NA-AR(4) & 1.36 [1.02, 1.75] & \textbf{4.23} [3.53, 4.85] \\
OS scale, $k=8$ & 1.22 [0.98, 1.50] & 1.17 [0.80, 1.53] \\
Power OS & 1.03 [1.01, 1.06] & 1.09 [0.91, 1.28] \\
P-ANMF, maximum & 1.16 [0.85, 1.61] & \textbf{4.33} [3.67, 5.15] \\
P-ANMF, one bin & 1.15 [0.79, 1.75] & \textbf{3.50} [2.77, 4.24] \\
PAMF-H, one Doppler & 1.15 [0.98, 1.42] & \textbf{4.09} [3.32, 5.08] \\
IN-ARCP, residual bootstrap & --- & 0.75 [0.32, 1.16] \\
\midrule\multicolumn{3}{c}{$\alpha=10^{-4}$}\\
ANMF-Tyler & 1.18 [0.91, 1.63] & \textbf{4.73} [4.08, 5.81] \\
ANMF-SCM & 1.10 [0.91, 1.41] & \textbf{5.64} [4.90, 6.63] \\
Power CA & 1.18 [1.04, 1.33] & \textbf{0.17} [0.07, 0.25] \\
Clipped, OS scale & \textbf{0.39} [0.13, 0.65] & \textbf{11.86} [8.29, 16.54] \\
Integration, AR(4) & \textbf{4.26} [1.46, 9.24] & \textbf{7.78} [3.95, 13.54] \\
IN-ARCP & 1.30 [0.94, 1.72] & 1.70 [0.91, 2.66] \\
IN-AR(4) & 1.31 [1.02, 1.60] & 1.85 [1.05, 3.07] \\
NA-AR(4) & \textbf{2.08} [1.14, 3.34] & \textbf{12.69} [9.13, 16.26] \\
OS scale, $k=8$ & 1.38 [0.82, 2.05] & \textbf{2.76} [1.25, 4.85] \\
Power OS & 1.02 [0.91, 1.10] & \textbf{2.21} [1.69, 2.78] \\
P-ANMF, maximum & \textbf{2.33} [0.85, 5.19] & \textbf{9.25} [6.27, 13.24] \\
P-ANMF, one bin & \textbf{2.52} [0.80, 6.10] & \textbf{6.59} [3.88, 11.12] \\
PAMF-H, one Doppler & \textbf{2.18} [1.18, 3.78] & \textbf{8.49} [5.83, 12.23] \\
IN-ARCP, residual bootstrap & --- & 1.22 [0.52, 1.93] \\
\bottomrule
\end{tabular}
\end{table}

\subsection{Pivotality across clutter power, all thirteen statistics}\label{sup:pivot}
Table~\ref{tab:pivot} gives the texture-quintile spread, the ratio of the largest to the
smallest measured false-alarm rate over quintiles of local clutter power, for every statistic
of Table~\ref{M-tab:lowpfa} under both threshold rules. The spread is a pre-specified outcome
of the frozen protocol; the two readings of it below are exploratory.

\emph{The remedy recommended for persistent targets has a strongly texture-dependent rate.}
Proposition~\ref{M-prop:blind} rules out history normalization for a target present throughout
the history, and the paper recommends the self-normalized whitened Doppler statistic instead.
Its conformal rate is compliant when pooled (1.04 at $10^{-2}$, 1.47 at $10^{-4}$) but varies
by a factor of 23 across texture quintiles at $10^{-4}$, against 2.7 for
IN-ARCP. It carries the widest spread in the table at $10^{-2}$, though not at the two lower
design rates, where the saturating clipped integrator is wider still. The statistic is scale-invariant and therefore texture-pivotal in pure
compound-Gaussian clutter. Additive noise can break this invariance
(Corollary~\ref{M-cor:texture}); residual correlation and model mismatch can also contribute.
The quintile spreads do not isolate the cause. The dwell stays in one range cell.

\emph{The model-based laws fail in two distinct ways.} Among the rules that exceed twice design
at $10^{-4}$, the order-statistic and maximum-DFT laws carry the largest texture spreads
(4.0 and 4.5 at $10^{-3}$, against 1.4 for the CA law of IN-ARCP), so their
rate depends strongly on clutter power. The summary prints only the max/min ratio, but the saved outcome arrays retain all five
quintile exceedance counts and denominators. No direction or cause follows from a spread
alone. Three other rules behave differently:
the clipped integrator's simulated threshold fails worse still in the pooled rate (17.19 at
$10^{-3}$, 59.89 at $10^{-4}$) and the two ANMF laws less badly (3.57 and 2.66 at $10^{-3}$),
but all three almost evenly across quintiles (1.3, 2.0 and 2.2). For the
Kraut--Scharf law that is expected, since the paper's simulation shows it already failing
2.8--3.4-fold at $10^{-2}$ in compound-Gaussian clutter, where the real failure is smaller.
Attributing the difference to tail weight would require an analysis beyond these spread
summaries; it is not established here.

\begin{table}[h]
\centering\small
\caption{Texture-quintile spread of the measured false-alarm rate (max/min over quintiles of
local clutter power, $\pm1024$-pulse proxy), IPIX, 56 units, from \texttt{r12/analyze\_r12.py}
Part L. Conformal: split-conformal threshold. Model: the rule named in
Table~\ref{M-tab:lowpfa}. The clipped integrator's conformal spread at $10^{-4}$ is unbounded
because some quintiles record no false alarm at all, its statistic having saturated.}
\label{tab:pivot}
\begin{tabular}{lcccccc}
\toprule
 & \multicolumn{3}{c}{Conformal} & \multicolumn{3}{c}{Model-based rule}\\
\cmidrule(lr){2-4}\cmidrule(lr){5-7}
$\alpha$ & $10^{-2}$ & $10^{-3}$ & $10^{-4}$ & $10^{-2}$ & $10^{-3}$ & $10^{-4}$\\
\midrule
IN-ARCP (CA law) & 1.3 & 1.5 & 2.7 & 1.3 & 1.4 & 2.0\\
IN-AR(4) (CA law) & 1.1 & 1.1 & 1.9 & 1.1 & 1.1 & 1.3\\
OS scale, $k=8$ (OS law) & 1.7 & 1.9 & 3.0 & 2.1 & 4.0 & 7.1\\
NA-AR(4) (Gaussian) & 1.9 & 3.0 & 2.9 & 1.6 & 2.3 & 3.2\\
Power, CA (white law) & 1.1 & 1.1 & 1.3 & 1.0 & 1.3 & 1.6\\
Power, OS (white law) & 1.1 & 1.9 & 3.3 & 1.1 & 1.4 & 1.8\\
Integration (dwell law) & 1.7 & 1.8 & 2.2 & 1.6 & 2.1 & 2.0\\
PAMF-H, one Doppler (CA law) & 2.2 & 2.1 & 4.0 & 2.2 & 2.4 & 2.4\\
P-ANMF, one bin (Beta) & 1.5 & 1.7 & 2.0 & 2.3 & 2.4 & 2.0\\
P-ANMF, max (Fisher's $g$) & 4.4 & 9.9 & 23 & 2.3 & 4.5 & 7.0\\
ANMF-SCM (Kraut--Scharf) & 2.3 & 1.8 & 3.2 & 2.2 & 2.0 & 1.8\\
ANMF-Tyler (FP law) & 2.6 & 2.1 & 2.0 & 2.5 & 2.2 & 1.8\\
Clipped, OS scale (Gaussian sim.) & 2.0 & 12 & $>10^{4}$ & 1.2 & 1.3 & 1.5\\
\bottomrule
\end{tabular}
\end{table}


\section{Coverage, width and transfer checks}
\subsection{Synthetic verification}
Figure~\ref{fig:sup_synthetic} compares exact and simulated clutter-to-noise-ratio-conditional coverage at nominal 0.90.
\begin{figure}[htbp]
\centering\includegraphics[width=7.16in]{fig_synthetic_exact.pdf}
\caption{Synthetic verification: coverage conditional on the true clutter-to-noise ratio. The blue line is the exact law of Proposition~\ref{M-prop:law}; open circles are its simulation. Other markers show the comparison methods. Each panel uses 200 replications with $m=16$ at nominal coverage 0.90; Monte Carlo standard error is at most 0.0014.}
\label{fig:sup_synthetic}
\end{figure}


\subsection{Texture-conditional coverage and disc widths on IPIX}
Figure~\ref{fig:sup_texture} compares texture-conditional coverage; the following table reports coverage and disc widths.
\begin{figure}[htbp]
\centering\includegraphics[width=7.16in]{fig_texture_conditional.pdf}
\caption{Coverage by texture quintile on IPIX (rotations 2--3, $m=16$). Texture is proxied by the local clutter power ($\pm1024$ pulses, excluding the episode). Shading shows the expected maximum quintile deviation from sampling noise alone; it is a reference band, not a confidence interval for the plotted curves.}
\label{fig:sup_texture}
\end{figure}
\begin{table}[h]
\centering\small
\caption{IPIX width and coverage ($m=16$; rotations 2--3; 56 units). Radius ratio: geometric mean over units of mean disc radius relative to IN-ARCP on the same test episodes, with 95\% day-cluster bootstrap interval. Max.\ dev.: mean over units of the largest absolute texture-quintile coverage deviation from nominal. At equal model order, NA-AR(4) has 0.891 [0.872, 0.905] times the radius of the innovation-normalized AR(4) disc.}
\resizebox{\textwidth}{!}{\input{tab_main.tex}}
\end{table}

\subsection{Day transfer and texture-proxy sensitivity}
\begin{table}[h]
\centering\small
\caption{Day transfer ($m=16$, $1-\alpha=0.90$): fit and calibrate on other days, test on each of six held-out days.}
\input{tab_transfer.tex}
\end{table}
\noindent Texture-proxy sensitivity (pre-specified windows; mean maximum quintile deviation at $1-\alpha=0.90$, $m=16$): with $\pm512$ pulses IN-ARCP 0.035, NA-AR(4) 0.043, unnormalized 0.181, neural center 0.103; with $\pm2048$ pulses IN-ARCP 0.033, unnormalized 0.141. The ranking is the same as with $\pm1024$.


\subsection{Hold-out series}
Two further McMaster series (\#269 and \#287, the widely used ``high'' and ``low'' sea-state files, one provider-preprocessed VV cell each, from the same 1993 campaign) were hashed and left unopened until the main analysis was complete, and then evaluated once (Table~\ref{tab:holdout}). They confirm the direction of the main result, with larger gains. They also show a case in which the classical Gaussian threshold over-covers while conformal IN-ARCP stays near nominal. Their test sets are small (about 165 episodes per unit), so the hold-out is a directional check. Mondrian calibration returns infinite regions at $1-\alpha=0.99$ there, because each bin receives fewer calibration episodes than the rank requires.
\begin{table}[h]
\centering\small
\caption{Hold-out series \#269 and \#287 ($1-\alpha=0.90$; six units per column).}
\label{tab:holdout}
\input{tab_holdout.tex}
\end{table}


\subsection{Dependence, day sensitivity and optimizer diagnostics}
At $\alpha=0.1$, mean lag-one exceedance autocorrelation within an IPIX bin (256-pulse separation) is $-0.007$, with unit range $[-0.067,0.046]$; same-time cross-bin correlation is 0.045. These diagnostics do not prove independence. Observed maximum quintile coverage deviations for IN-ARCP, IN-AR(4), NA-AR(4) and the unnormalized score are 0.036, 0.032, 0.043 and 0.165. The binomial reference is 0.029; dependence-aware time-shift references are 0.033, 0.033, 0.039 and 0.062.

The AR(1) clamp binds in 5 of 56 IPIX units. NA-AR(2) converges in every unit; NA-AR(4) converges in 16 at the 1,500-iteration cap and 55 after a further 4,500 iterations. The exploratory restart changes its radius by geometric factors 0.9989 at coverage 0.90 (unit range 0.9790--1.0132) and 0.9997 at 0.99 (0.9492--1.0327); coverage changes from 0.9019 to 0.9020 and from 0.9905 to 0.9906. Median likelihood gain per training episode is $3.13\times10^{-6}$, maximum $1.54\times10^{-4}$. The median whole-unit GPU time for thirteen procedures is 71~s, not a detector hardware benchmark.


\begin{table}[htbp]
\centering\small\setlength{\tabcolsep}{4pt}
\caption{Leave-one-day-out ranges from \texttt{r10/audit\_existing.txt}. Gains are dB; radius ratio and spread are dimensionless. These complement the approximate six-day cluster intervals.}\label{tab:lodo}
\begin{tabular}{lcc}
\toprule
Effect & All six days & LODO range \\
\midrule
NA-AR(4)/IN-AR(4) radius ratio (0.90) & 0.891 & 0.884--0.895 \\
IN1 quintile Pfa spread ($\alpha=0.01$) & 1.75$\times$ & 1.72--1.83x \\
single-pulse gain IN1 vs CA16  $P_{\rm d}=0.5$ & 6.96 dB & 5.97-- 8.16 \\
single-pulse gain IN1 vs CAloc $P_{\rm d}=0.5$ & 9.81 dB & 8.94--11.54 \\
single-pulse gain NA4 vs IN1   $P_{\rm d}=0.5$ & 1.07 dB & 1.03-- 1.11 \\
single-pulse gain NA4 vs IN4   $P_{\rm d}=0.5$ & 0.66 dB & 0.58-- 0.75 \\
single-pulse gain IN1 vs CA16  $P_{\rm d}=0.9$ & 7.62 dB & 6.59-- 9.01 \\
single-pulse gain IN1 vs CAloc $P_{\rm d}=0.9$ & 8.97 dB & 8.10--10.62 \\
single-pulse gain NA4 vs IN1   $P_{\rm d}=0.9$ & 2.19 dB & 1.88-- 2.39 \\
single-pulse gain NA4 vs IN4   $P_{\rm d}=0.9$ & 1.46 dB & 1.43-- 1.52 \\
eight-look dwell gain IN1 vs CAloc-sum $P_{\rm d}=0.5$ & 10.70 dB & 9.86--12.38 \\
eight-look dwell gain NA4 vs CAloc-sum $P_{\rm d}=0.5$ & 11.74 dB & 11.01--13.00 \\
\bottomrule
\end{tabular}
\end{table}

The six daily onset gains of IN-ARCP over local power at $P_{\rm d}=0.5$ are 3.5, 8.4, 13.8, 13.3, 12.2 and 11.4~dB. They should not be read as interchangeable independent replications.

\subsection{Parametric texture and adaptive conformal baselines}
A fitted K-texture law changes NA-AR(4)'s radius by factors 1.0092 [1.0064, 1.0115] at coverage 0.90 and 1.0078 [1.0034, 1.0104] at 0.99; coverage and quintile deviations are essentially unchanged. Within-session ACI with $\gamma=0.005$ or 0.02 gives similar coverage and finite radii; at 0.99, $\gamma=0.02$ gives infinite regions in 0.0115 of episodes. Day-transfer worst coverages improve from 0.8822 for split IN-ARCP to 0.8889 and 0.8965 at coverage 0.90, and from 0.9875 to 0.9889 and 0.9902 at 0.99. Clean-feedback behavior does not establish safety under target-contaminated online feedback: at 5\% contamination ungated ACI gives $0.14\alpha$ with $P_{\rm d}=0.275$ (split: 0.890); gated ACI still gives only $0.55\alpha$ or $0.68\alpha$. Full settings, trajectories and uncertainty remain in the baseline and R10 outputs.

\section{Detection comparisons and target-history contamination}\label{sup:detect}
Injected comparisons share known timing and target-aligned dwells. They test mechanisms, not operational acquisition. Figure~\ref{fig:sup_detection} retains the $10^{-3}$ design rate and full method curves.
\begin{figure}[htbp]\centering\includegraphics[width=7.16in]{fig_detection.pdf}
\caption{Injected Swerling-1 onset detection on IPIX at design false-alarm probabilities 0.01 and 0.001 ($m=16$, conformal thresholds). Empirical and law curves are unweighted unit means; circles show the fitted-model law of Corollary~\ref{M-cor:detect}. Reported SCR gains average per-unit gains.}\label{fig:sup_detection}
\end{figure}


\begin{table}[htbp]
\centering\small\setlength{\tabcolsep}{4pt}
\caption{Onset SCR gain at design $P_{\rm fa}=0.01$: mean per-unit gain [95\% day-cluster interval], from \texttt{detection/detection\_summary.json}. Positive means less SCR. Local power uses the non-causal $\pm1024$-pulse window.}\label{tab:onset_gains}
\begin{tabular}{lccc}
\toprule
Method & $P_{\rm d}$ & Gain over CA16 (dB) & Gain over local power (dB) \\
\midrule
IN1 & $0.5$ & 6.96 [4.29, 9.04] & 9.81 [6.41, 12.61] \\
IN4 & $0.5$ & 7.37 [4.54, 9.52] & 10.22 [6.67, 13.26] \\
NA4 & $0.5$ & 8.03 [5.45, 10.02] & 10.88 [7.56, 13.66] \\
G1 & $0.5$ & 7.05 [4.40, 9.11] & 9.90 [6.48, 12.71] \\
U & $0.5$ & 2.30 [0.78, 3.44] & 5.15 [3.34, 6.61] \\
MLP & $0.5$ & 3.88 [1.89, 5.34] & 6.73 [3.89, 8.70] \\
LS & $0.5$ & 5.72 [2.75, 7.96] & 8.57 [4.89, 11.49] \\
IN1 & $0.9$ & 7.62 [4.55, 10.01] & 8.97 [5.58, 11.67] \\
IN4 & $0.9$ & 8.35 [4.95, 10.96] & 9.70 [5.90, 12.64] \\
NA4 & $0.9$ & 9.81 [6.27, 12.45] & 11.16 [7.28, 14.14] \\
G1 & $0.9$ & 7.67 [4.71, 10.02] & 9.02 [5.70, 11.73] \\
U & $0.9$ & 2.80 [1.01, 4.34] & 4.15 [2.23, 5.73] \\
MLP & $0.9$ & 4.36 [2.30, 5.85] & 5.71 [3.23, 7.35] \\
LS & $0.9$ & 6.14 [2.75, 8.69] & 7.49 [3.78, 10.46] \\
\bottomrule
\end{tabular}
\end{table}

\section{Dwell detection on IPIX}\label{sup:dwell}
\begin{table}[h]
\centering\small
\caption{SCR (dB) for $P_{\rm d}=0.5$ and $0.9$ over an 8-pulse dwell: persistent Swerling-1 target with random Doppler from the first dwell pulse (IPIX, conformal thresholds). $\le-5$: reached at the lowest SCR tested; n.r.: not reached at 25~dB.}
\label{tab:dwellsup}
\input{tab_dwell.tex}
\end{table}

PAMF-H has the best tested dwell sensitivity. Its endpoint at the lowest tested SCR makes the ANMF-Tyler gap of 5.36 [1.93, 8.83]~dB a lower bound. Clipping costs 2.04 [0.00, 2.90]~dB against AR(1) integration; the onset raw-power-OS/IN-ARCP gap is 10.62 [7.05, 12.69]~dB. The earlier R10 equal-dwell comparison gives PAMF-H a 1.0 [0.3, 1.7]~dB advantage over integrated IN-ARCP at $P_{\rm d}=0.5$ (about 4~dB at 0.9); the studies have separate calibrations. A Hann-windowed range-CA Doppler bank needs 5.0 [4.2, 6.0]~dB more SCR than integrated IN-ARCP. PAMF-H's matched-Doppler cost in that study is 4.25~dB.

\section{Guarded scale, clip level, certified trimming, pulse blanking and bursts (R14)}\label{sup:r14}
Frozen protocol \texttt{r14/PROTOCOL\_R14.md} (sha256 9f771091\ldots), written after a synthetic dry run and before any IPIX outcome; expectations and verdicts in \texttt{r14/RESULTS\_R14.md} (12 of 15 met). The settings are those of the final IPIX study, except for dense calibration: non-overlapping 24-pulse segments, about 14,600 per unit, so that $\alpha=10^{-3}$ is resolved. The study tested five changes:
\begin{itemize}
\item a guard of $\Delta\in\{0,8,16\}$ pulses between the scale window and the tested pulse;
\item clip levels $\kappa\in\{4,6,9,12,18\}$;
\item a certified trimmed sum, which discards the largest dwell terms; the number discarded is the smallest for which more corrupted dwell innovations than discarded terms have probability at most $\alpha/5$ under the hit model (4 of eight at a 2\% and 5 at a 5\% hit rate at $\alpha=10^{-2}$; no such number exists for bursts);
\item pulse blanking of isolated spikes, calibrated with masks drawn from the hit model and passed through the same rule (mask-matched calibration);
\item bursts of eight consecutive pulses at the same mean hit rate.
\end{itemize}
The certified thresholds use the worst case of Theorem~\ref{M-thm:cert} with hits drawn either from the independent-hit model or from the burst model. The guard leaves the false-alarm law unchanged; the measured $P_{\rm fa}/\alpha$ with $\Delta=8$ is 1.03 at $10^{-2}$ and 1.03 at $10^{-3}$. The real IPIX target stays undetected with every guard (exceedance 0.0005 at $10^{-3}$ with $\Delta=8$), because it is present in every history.

\emph{Unmet expectations of R14.}
\begin{itemize}
\item \textbf{Onset cost (G3).} The guard was expected to cost at most 0.5~dB at onset. It costs 0.9~dB with $\Delta=8$ and 1.2~dB with $\Delta=16$, because the scale is older.
\item \textbf{Certification at $10^{-3}$ (K2).} The certificate at the clip level chosen by the saturation rule was expected to reach $P_{\rm d}=0.5$ below 25~dB in all four interference conditions at $10^{-3}$. With 5\% hits the certificate is vacuous at $10^{-3}$ for every clip level.
\item \textbf{Cost of trimming (T2).} Trimming was expected to cost more than clipping at both hit rates. At 5\% it costs +5.1~dB against +5.9~dB for clipping ($\kappa=6$).
\end{itemize}
The SCR for $P_{\rm d}=0.5$ is fragile under 5\% hits, because the certified detection curve is nearly flat there, so the main text reports $P_{\rm d}$ at 10~dB.



\begin{table}[htbp]
\centering\small\setlength{\tabcolsep}{4pt}
\caption{Guarded scale: clean rate ratio [cluster interval], and mean per-look $P_{\rm d}$ over looks 1--8 at 10~dB. R14 IPIX and R15 NetRAD; these eight looks do not measure post-guard persistence.}\label{tab:guards}
\begin{tabular}{lcccccc}
\toprule
Radar & $\Delta$ & $\alpha=10^{-2}$ & $10^{-3}$ & Random & Opposite & Matched \\
\midrule
IPIX & 0 & 1.01 [0.98, 1.04] & 1.02 [0.94, 1.10] & 0.24 & 0.36 & 0.01 \\
IPIX & 8 & 1.03 [0.97, 1.07] & 1.03 [0.84, 1.21] & 0.82 & 0.95 & 0.07 \\
IPIX & 16 & 1.00 [0.93, 1.03] & 1.03 [0.81, 1.15] & 0.82 & 0.94 & 0.07 \\
NetRAD & 0 & 1.05 [0.98, 1.14] & 1.22 [0.94, 1.65] & 0.30 & 0.43 & 0.00 \\
NetRAD & 8 & 1.05 [1.00, 1.11] & 1.23 [0.98, 1.64] & 0.92 & 0.99 & 0.03 \\
NetRAD & 16 & 1.05 [1.00, 1.12] & 1.36 [1.05, 1.87] & 0.92 & 0.99 & 0.03 \\
\bottomrule
\end{tabular}
\end{table}


\begin{table}[htbp]
\centering\small\setlength{\tabcolsep}{4pt}
\caption{Detection within eight looks for linear emergence, 10~dB SCR and design $P_{\rm fa}=0.01$.}\label{tab:ramps}
\begin{tabular}{lcccc}
\toprule
Radar & Ramp length & $\Delta=0$ & $8$ & $16$ \\
\midrule
IPIX & 8 & 0.745 & 0.908 & 0.909 \\
IPIX & 16 & 0.032 & 0.869 & 0.903 \\
NetRAD & 8 & 0.862 & 0.964 & 0.966 \\
NetRAD & 16 & 0.016 & 0.937 & 0.959 \\
\bottomrule
\end{tabular}
\end{table}

The earlier pre-onset ramp comparison, corrected R10 verdict, and all R11 failed expectations remain in Section~\ref{sup:unmet}. Exact ramp-law error is 0.016, within the 0.03 expectation. The exploratory post-onset law has mean per-unit/look errors 0.010, 0.003 and 0.008 for random, matched and opposite Doppler. Cumulative-null and equal-dwell comparisons remain in the full onset outputs. The R11 median OS score has onset cost 1.99 [0.89, 3.39]~dB, quintile spread 2.06 and coverage-law error 0.0214; these miss the stated tolerances.


\begin{table}[htbp]
\centering\small\setlength{\tabcolsep}{4pt}
\caption{Clean clipped integration, IPIX R14 dense calibration: rate ratio [95\% day-cluster interval].}\label{tab:clip_scan}
\begin{tabular}{lcc}
\toprule
$\kappa$ & $\alpha=10^{-2}$ & $10^{-3}$ \\
\midrule
4 & 1.04 [0.97, 1.08] & 0.09 [0.00, 0.23] \\
6 & 1.05 [0.99, 1.08] & 0.67 [0.35, 0.94] \\
9 & 1.05 [0.98, 1.09] & 0.90 [0.58, 1.15] \\
12 & 1.06 [1.00, 1.09] & 0.95 [0.65, 1.20] \\
18 & 1.06 [1.01, 1.09] & 1.03 [0.79, 1.19] \\
\bottomrule
\end{tabular}
\end{table}

The pre-specified saturation rule chooses $\kappa=9$ in 17 of 56 units at $10^{-3}$ and finds no eligible tested level ($\kappa\le18$) in 16. Applying $\kappa=9$ to every unit is post hoc: at $10^{-3}$ it changes 10-dB detection from 0.46 to 0.68. Sparse R12 and dense R14 calibrations are separate analyses; their endpoints must not be combined.


\begin{table}[htbp]
\centering\small\setlength{\tabcolsep}{4pt}
\caption{IPIX independent-hit remedies at $\alpha=10^{-2}$: rate ratio [95\% day-cluster interval], detection at 10 and 25~dB SCR. R14 dense calibration; interference is 10 or 30~dB above local power. Certificate calibration uses matching independent hits.}\label{tab:remedies_2}
\begin{tabular}{lccccc}
\toprule
Hit rate & JNR (dB) & Statistic & Rate ratio & $P_{\rm d}(10)$ & $P_{\rm d}(25)$ \\
\midrule
2\% & 10 & Clean clip6 & 1.08 [1.05, 1.11] & 0.85 & 0.99 \\
2\% & 10 & Certified clip6 & 0.53 [0.43, 0.61] & 0.81 & 0.95 \\
2\% & 10 & Certified clip9 & 0.52 [0.44, 0.58] & 0.81 & 0.96 \\
2\% & 10 & Certified trim4 & 0.53 [0.48, 0.57] & 0.79 & 0.97 \\
2\% & 10 & Blanking & 1.05 [1.01, 1.08] & 0.79 & 0.97 \\
2\% & 30 & Clean clip6 & 1.10 [1.07, 1.12] & 0.85 & 0.99 \\
2\% & 30 & Certified clip6 & 0.55 [0.46, 0.61] & 0.81 & 0.95 \\
2\% & 30 & Certified clip9 & 0.53 [0.45, 0.58] & 0.81 & 0.96 \\
2\% & 30 & Certified trim4 & 0.57 [0.51, 0.60] & 0.79 & 0.97 \\
2\% & 30 & Blanking & 1.08 [1.03, 1.11] & 0.79 & 0.97 \\
5\% & 10 & Clean clip6 & 1.22 [1.15, 1.31] & 0.85 & 0.98 \\
5\% & 10 & Certified clip6 & 0.12 [0.09, 0.14] & 0.64 & 0.75 \\
5\% & 10 & Certified clip9 & 0.13 [0.08, 0.17] & 0.71 & 0.87 \\
5\% & 10 & Certified trim5 & 0.15 [0.12, 0.17] & 0.72 & 0.95 \\
5\% & 10 & Blanking & 1.13 [1.04, 1.21] & 0.78 & 0.96 \\
5\% & 30 & Clean clip6 & 1.28 [1.21, 1.37] & 0.85 & 0.98 \\
5\% & 30 & Certified clip6 & 0.12 [0.10, 0.14] & 0.64 & 0.75 \\
5\% & 30 & Certified clip9 & 0.14 [0.10, 0.18] & 0.71 & 0.87 \\
5\% & 30 & Certified trim5 & 0.23 [0.19, 0.26] & 0.72 & 0.95 \\
5\% & 30 & Blanking & 1.21 [1.16, 1.25] & 0.78 & 0.96 \\
\bottomrule
\end{tabular}
\end{table}


\begin{table}[htbp]
\centering\small\setlength{\tabcolsep}{4pt}
\caption{IPIX independent-hit remedies at $\alpha=10^{-3}$: rate ratio [95\% day-cluster interval], detection at 10 and 25~dB SCR. R14 dense calibration; interference is 10 or 30~dB above local power. Certificate calibration uses matching independent hits.}\label{tab:remedies_3}
\begin{tabular}{lccccc}
\toprule
Hit rate & JNR (dB) & Statistic & Rate ratio & $P_{\rm d}(10)$ & $P_{\rm d}(25)$ \\
\midrule
2\% & 10 & Clean clip6 & 0.65 [0.34, 0.92] & 0.46 & 0.54 \\
2\% & 10 & Certified clip6 & 0.07 [0.03, 0.11] & 0.12 & 0.16 \\
2\% & 10 & Certified clip9 & 0.27 [0.17, 0.34] & 0.52 & 0.63 \\
2\% & 10 & Certified trim5 & 0.39 [0.30, 0.44] & 0.68 & 0.94 \\
2\% & 10 & Blanking & 1.09 [1.00, 1.17] & 0.68 & 0.93 \\
2\% & 30 & Clean clip6 & 0.65 [0.33, 0.90] & 0.46 & 0.54 \\
2\% & 30 & Certified clip6 & 0.07 [0.03, 0.11] & 0.12 & 0.16 \\
2\% & 30 & Certified clip9 & 0.26 [0.16, 0.33] & 0.52 & 0.63 \\
2\% & 30 & Certified trim5 & 0.43 [0.33, 0.49] & 0.68 & 0.94 \\
2\% & 30 & Blanking & 1.13 [1.04, 1.18] & 0.68 & 0.93 \\
5\% & 10 & Clean clip6 & 0.66 [0.35, 0.94] & 0.45 & 0.54 \\
5\% & 10 & Certified clip6 & 0.00 [0.00, 0.00] & 0.00 & 0.00 \\
5\% & 10 & Certified clip9 & 0.00 [0.00, 0.00] & 0.00 & 0.00 \\
5\% & 10 & Certified trim6 & 0.00 [0.00, 0.00] & 0.00 & 0.00 \\
5\% & 10 & Blanking & 1.67 [1.39, 1.91] & 0.67 & 0.93 \\
5\% & 30 & Clean clip6 & 0.69 [0.39, 0.93] & 0.45 & 0.54 \\
5\% & 30 & Certified clip6 & 0.00 [0.00, 0.00] & 0.00 & 0.00 \\
5\% & 30 & Certified clip9 & 0.00 [0.00, 0.00] & 0.00 & 0.00 \\
5\% & 30 & Certified trim6 & 0.00 [0.00, 0.00] & 0.00 & 0.00 \\
5\% & 30 & Blanking & 2.25 [1.87, 2.54] & 0.67 & 0.93 \\
\bottomrule
\end{tabular}
\end{table}


\begin{table}[htbp]
\centering\small\setlength{\tabcolsep}{4pt}
\caption{Eight-pulse bursts on IPIX at JNR 30~dB: both hit rates and both design rates, including burst-matched calibration. Burst certificates also detect no target at 25~dB.}\label{tab:bursts}
\begin{tabular}{lcccc}
\toprule
Hit rate & $\alpha$ & Statistic & Rate ratio [95\% CI] & $P_{\rm d}(10)$ \\
\midrule
2\% & $10^{-2}$ & Clean clip6 & 2.24 [2.11, 2.33] & 0.84 \\
2\% & $10^{-2}$ & Independent-hit cert., clip6 & 1.46 [1.33, 1.58] & 0.80 \\
2\% & $10^{-2}$ & Independent-hit cert., clip9 & 1.75 [1.59, 1.89] & 0.80 \\
2\% & $10^{-2}$ & Burst certificate, clip6 & 0.00 [0.00, 0.00] & 0.00 \\
2\% & $10^{-2}$ & Blanking, independent masks & 2.31 [2.22, 2.36] & 0.79 \\
2\% & $10^{-2}$ & Blanking, burst masks & 2.25 [2.17, 2.31] & 0.79 \\
2\% & $10^{-3}$ & Clean clip6 & 3.86 [2.11, 5.34] & 0.45 \\
2\% & $10^{-3}$ & Independent-hit cert., clip6 & 0.90 [0.38, 1.61] & 0.12 \\
2\% & $10^{-3}$ & Independent-hit cert., clip9 & 3.86 [2.78, 4.99] & 0.51 \\
2\% & $10^{-3}$ & Burst certificate, clip6 & 0.00 [0.00, 0.00] & 0.00 \\
2\% & $10^{-3}$ & Blanking, independent masks & 13.60 [12.63, 14.17] & 0.69 \\
2\% & $10^{-3}$ & Blanking, burst masks & 13.43 [12.42, 14.00] & 0.68 \\
5\% & $10^{-2}$ & Clean clip6 & 3.98 [3.65, 4.22] & 0.82 \\
5\% & $10^{-2}$ & Independent-hit cert., clip6 & 1.25 [1.04, 1.45] & 0.62 \\
5\% & $10^{-2}$ & Independent-hit cert., clip9 & 1.66 [1.44, 1.83] & 0.69 \\
5\% & $10^{-2}$ & Burst certificate, clip6 & 0.00 [0.00, 0.00] & 0.00 \\
5\% & $10^{-2}$ & Blanking, independent masks & 4.14 [4.01, 4.23] & 0.78 \\
5\% & $10^{-2}$ & Blanking, burst masks & 4.10 [3.98, 4.17] & 0.78 \\
5\% & $10^{-3}$ & Clean clip6 & 8.92 [5.10, 12.32] & 0.44 \\
5\% & $10^{-3}$ & Independent-hit cert., clip6 & 0.00 [0.00, 0.00] & 0.00 \\
5\% & $10^{-3}$ & Independent-hit cert., clip9 & 0.00 [0.00, 0.00] & 0.00 \\
5\% & $10^{-3}$ & Burst certificate, clip6 & 0.00 [0.00, 0.00] & 0.00 \\
5\% & $10^{-3}$ & Blanking, independent masks & 32.43 [30.34, 33.64] & 0.68 \\
5\% & $10^{-3}$ & Blanking, burst masks & 32.35 [30.32, 33.49] & 0.67 \\
\bottomrule
\end{tabular}
\end{table}

At $10^{-3}$ and 5\% independent hits, every tested clip certificate is vacuous. At $10^{-2}$ and 5\% hits, trimming costs 5.1~dB versus clipping's 5.9~dB for $\kappa=6$, contrary to the expectation that trimming costs more. Fixed-SCR detection is clearer because the certified curves are nearly flat. Burst-matched trimming has no usable choice. The R12 sparse-calibration binary certificate has 10-dB detection within 0.03 of certified clipping at both tested rates and powers. At design $\alpha=10^{-2}$, uncertified integration and PAMF-H exceed design 10--23-fold, while self-normalized statistics stay within 1.21 times design at a sensitivity cost. Full clip grids, SCR costs and sparse-calibration rows remain in the frozen outputs.

\subsection{Pulse blanking: the thresholds}\label{sup:blank}
Pulse blanking discards an innovation
power, and the innovation after it, when the power exceeds both ten times the history median
and four times a robust level of its block. That level is the history median for history terms
and the third-smallest dwell power for dwell terms. The remaining dwell powers are integrated.
Blanking is calibrated on clean episodes blanked at hit positions drawn from the hit model, so
it has a conformal threshold but no certificate: nothing bounds the statistic when a hit
survives the blanking rule.

\section{77~GHz thermal-noise and real-interference checks}\label{sup:jkusel}
Range bins 8--247 and receivers 0, 5, 10 and 15 of both stations are used. Figure~\ref{fig:sup_ground} shows coverage and onset gain by CNR class; the general-form gain curve is post hoc. IN-ARCP's gain over local power rises from $-1.45$~dB in noise-dominated cells to 3.11~dB at 10--20~dB CNR. Only 338 test episodes from two units occupy the highest CNR class, versus 721,095 in the lowest. At design 0.01 its CNR-class rate ratios are 1.07, 1.06, 0.95, 0.52, 0.08 and 0 across $<-5$, $[-5,0)$, $[0,5)$, $[5,10)$, $[10,20)$ and $\ge20$~dB; the support is uneven.
\begin{figure}[htbp]\centering\includegraphics[width=7.16in]{fig_jku.pdf}
\caption{77~GHz ground clutter: (a) coverage by CNR class at nominal 0.90; (b) onset SCR gain over local power for $P_{\rm d}=0.5$ at design $P_{\rm fa}=0.01$. The general-form gain interpretation is post hoc. CNR classes have unequal support, detailed in the text.}\label{fig:sup_ground}
\end{figure}

For real interference, calibration masks are random 24-chirp blocks from frames 0--49; tests use frames 50--99. N: no fast-time mitigation; Z: zeroing; ZAR: zeroing with AR(8) reconstruction. Test hit rates are 45.4\% in dense A and 4.6\% in sparse C. The empirical masks do not establish exchangeability, clutter independence or inclusion dominance. In sparse C without mitigation, integration runs at 3.56 times design in the 1--2-hit stratum. The clipped certificate costs 3.1~dB on hit-free segments and 6.4~dB overall against clean-threshold clipping. In dense A it never reaches $P_{\rm d}=0.5$ overall. AR reconstruction improves IN-ARCP SCR50 by only 0.18~dB, below the expected 1~dB.


\begin{table}[htbp]
\centering\small\setlength{\tabcolsep}{4pt}
\caption{Real 77~GHz interference, $\alpha=0.01$: all three mitigations. The last column is the hit-free stratum, not a confidence interval; n.r.: not reached in the tested SCR range. Full per-hit-stratum counts and ratios remain in \texttt{r12/r12\_jku\_summary.txt}.}\label{tab:jku_interference}
\begin{tabular}{lccccc}
\toprule
Run & Mitigation & Statistic & Rate ratio & SCR50 all & SCR50 clean \\
\midrule
ref & N & Integration & 1.03 & 5.91 & 5.91 \\
ref & N & Clean clip & 1.00 & 6.84 & 6.84 \\
ref & N & Certified clip & 1.00 & 6.84 & 6.84 \\
ref & N & Certified binary & 0.67 & 7.96 & 7.96 \\
ref & Z & Integration & 0.99 & 6.24 & 6.24 \\
ref & Z & Clean clip & 0.97 & 7.11 & 7.11 \\
ref & Z & Certified clip & 0.97 & 7.11 & 7.11 \\
ref & Z & Certified binary & 0.63 & 8.32 & 8.32 \\
ref & ZAR & Integration & 1.01 & 6.09 & 6.09 \\
ref & ZAR & Clean clip & 0.99 & 7.00 & 7.00 \\
ref & ZAR & Certified clip & 0.99 & 7.00 & 7.00 \\
ref & ZAR & Certified binary & 0.66 & 8.17 & 8.17 \\
A & N & Integration & 0.81 & 11.14 & 7.68 \\
A & N & Clean clip & 1.01 & 11.95 & 8.66 \\
A & N & Certified clip & 0.11 & n.r. & 13.74 \\
A & N & Certified binary & 0.06 & n.r. & 14.93 \\
A & Z & Integration & 0.96 & 6.96 & 6.11 \\
A & Z & Clean clip & 0.98 & 7.86 & 7.05 \\
A & Z & Certified clip & 0.11 & n.r. & 12.24 \\
A & Z & Certified binary & 0.06 & n.r. & 13.63 \\
A & ZAR & Integration & 1.01 & 6.74 & 6.06 \\
A & ZAR & Clean clip & 1.01 & 7.65 & 7.03 \\
A & ZAR & Certified clip & 0.11 & n.r. & 12.18 \\
A & ZAR & Certified binary & 0.06 & n.r. & 13.52 \\
C & N & Integration & 1.42 & 7.15 & 6.34 \\
C & N & Clean clip & 1.29 & 8.11 & 7.29 \\
C & N & Certified clip & 0.54 & 14.52 & 10.40 \\
C & N & Certified binary & 0.33 & 15.78 & 11.63 \\
C & Z & Integration & 1.00 & 6.41 & 6.22 \\
C & Z & Clean clip & 1.02 & 7.26 & 7.08 \\
C & Z & Certified clip & 0.51 & 13.94 & 10.00 \\
C & Z & Certified binary & 0.27 & 15.19 & 11.33 \\
C & ZAR & Integration & 1.04 & 6.25 & 6.10 \\
C & ZAR & Clean clip & 1.02 & 7.15 & 7.01 \\
C & ZAR & Certified clip & 0.50 & 13.88 & 9.93 \\
C & ZAR & Certified binary & 0.30 & 15.08 & 11.22 \\
\bottomrule
\end{tabular}
\end{table}

\section{The real IPIX target: detection versus dwell length}\label{sup:target}
\begin{table}[h]
\centering\small
\caption{Real IPIX target (primary cell, test thirds, 56 units): mean $P_{\rm d}$ at design $P_{\rm fa}=10^{-3}$ versus dwell $N$. In parentheses: measured clutter $P_{\rm fa}$ ($\times10^{-3}$). P-ANMF with Fisher's $g$ uses the analytic threshold.}
\label{tab:target}
\input{tab_target.tex}
\end{table}
Calibration of the self-normalized detector degrades for longer, overlapping windows. Its analytic Fisher's $g$ threshold fails already at $N=8$, consistent with imperfect whitening by an AR(4) model fitted on neighboring cells. Conformal P-ANMF retains sensitivity to the persistent target without history self-masking; its calibration degrades at the longest tested windows.


History-normalized IN-ARCP and clipping exceed on the primary target cell at only 0.0004 at design $10^{-3}$. This supports full-history self-masking, not real-onset validation. The long-dwell P-ANMF rate exceeds twice design; no i.i.d. Beta guarantee follows for the overlapping real-data windows. The one-second comparison with target-trained literature is descriptive because training and evaluation protocols differ.

\section{Second sea-clutter radar: NetRAD (R15)}\label{sup:r15}
Frozen protocol \texttt{r15/PROTOCOL\_R15.md} (sha256 945c5e51\ldots). It was written before any statistic of the NetRAD samples was computed. Before that point only the following had been read: file sizes, array shapes, and the providers' scripts and log. In addition, two of the providers' own quick-look images were viewed while the dataset was being located. Expectations and verdicts are in \texttt{r15/RESULTS\_R15.md} (12 of 15 met), and deviations in \texttt{r15/DEVIATIONS\_R15.md}. The deviations are a corrected site description (the Atlantic side of the Cape Peninsula, not False Bay) and run mechanics; none changes a method, setting or expectation.

\emph{Data.} Node 3 (monostatic) of the NetRAD system: seven HH and seven VV recordings of 130{,}000 pulses at 1~kHz. The clutter cells are the providers' per-recording range windows (\texttt{NR\_Range\_Selection\_June9.m}). Each cell has its complex mean removed and is scaled to unit power. The cross-polarized recordings were set aside before any data were read. All methods are the frozen R12 and R14 code. The only change is that the ANMF secondary cells are restricted to those two to four cells away, which gives 15--30 vectors (20--35 on IPIX). The AR(1) clamp $|r|\le0.98$ binds in 12 of the 14 recordings.

\emph{Unmet expectations of R15.}
\begin{itemize}
\item \textbf{Calibration at $10^{-4}$ (N1).} At least 12 of 13 conformal ratios were expected within $[0.5,2]$ at $10^{-4}$; 7 were. Outside: Clip-OS1 (0.39), D-IN4 (4.26), NA4 (2.08), P-ANMF (2.33), P-ANMF(bin) (2.52), PAMF-H(w) (2.18). At $10^{-3}$ all 13 are within $[0.5,2]$, and the 12 non-saturating statistics lie within 1.01--1.36 times design.
\item \textbf{Pivotality ranking (N4).} IN-ARCP's quintile spread at $10^{-2}$ was expected to be at most 2 and below CA16's. It is 2.40 against 1.51. In the same analysis on IPIX it is 1.31 against 1.06, so this expectation had no support on IPIX either.
\item \textbf{Flagged pulses (W2).} Non-coherent integration with a clean threshold was expected to exceed $1.5\alpha$ in segments with flagged pulses. It stays at 1.08 times design.
\end{itemize}
Expectation N5 (PAMF-H within 0.5~dB of the best dwell detector) was met, but it is uninformative: every whitened dwell detector reaches $P_{\rm d}=0.5$ at the lowest SCR tested.

\emph{Exploratory (post hoc, after the outcomes and the claim audit).} Two scripts look at where the conformal exceedances at $10^{-4}$ fall: \texttt{r15/explore\_r15\_hotspots.py} and \texttt{r15/explore\_r15\_attribution.py}.
\begin{itemize}
\item \textbf{Integration and PAMF-H.} Most of their excess comes from two VV units (1450 and 1501, assignment 3). Their exceedances fall within one to three seconds of the test third and coincide with flagged pulses, while their calibration thirds contain no flagged pulse. Without these two units the pooled ratios are 1.56 (integration) and 1.33 (PAMF-H).
\item \textbf{P-ANMF statistics.} Their excess comes from one HH unit (1128, assignment 3) with frequent flagged pulses in both thirds. Fewer than half of its exceedances coincide with flagged pulses.
\item \textbf{Unexplained.} The flagged pulses do not explain the excess of the noise-aware score (2.08). Nor do they explain an excess of integration of more than five times design in one unit that has no flagged pulse at all (1508, assignment 2).
\end{itemize}
\textbf{What the flags are is not established.} \texttt{r15/explore\_r15\_flag\_extent.py} counts, in the raw recordings, the range bins above ten times their median:
\begin{itemize}
\item a flagged HH pulse has a median of about 108 such bins, of 1024, and 85--91\% of them lie outside the clutter window;
\item unflagged HH pulses with an elevated window bin have about 80.
\end{itemize}
The flags are therefore the upper tail of a phenomenon that extends over range. They may be wideband interference: the providers remove WiFi pulses in their own processing of the bistatic nodes. They may equally be range-extended HH sea spikes. The main text therefore calls them flagged pulses, not interference.

In the clutter cells, uncertified non-coherent integration runs at 1.05 times design at $10^{-2}$ and 1.20 at $10^{-3}$ over all test segments of the 7 HH recordings, the only ones with at least 0.5\% flagged pulses (VV: at most 0.1\%). Segments without a flagged pulse give 1.04 and 1.17; segments with flagged pulses give at most 1.08 and 1.33, with wide intervals at $10^{-3}$ (\texttt{r15/r15\_summary.txt}, Part W). A certificate whose hit model must cover these frequent, bursty masks is nearly vacuous ($P_{\rm d}=0.04$). Blanking with mask-matched calibration keeps $P_{\rm d}=0.92$ at 0.57 times design.


The full pooled low-rate intervals remain in Table~\ref{tab:netrad_ci}; removing hotspot units is not used to replace the headline results. At $10^{-3}$ the flagged-pulse certificate is wholly vacuous, whereas mask-matched blanking has rate ratio 0.24 (upper recording-cluster bound 0.49) and $P_{\rm d}=0.83$. Flag origin, unexplained excesses and the weak burst remedy remain unresolved.

\section{Real onsets of a walking person at 77~GHz (R17)}\label{sup:r17}
Frozen protocol \texttt{r17/PROTOCOL\_R17.md}, committed and pushed before the data were obtained; verdicts in \texttt{r17/RESULTS\_R17.md}. In the pedestrian run of the 77~GHz data (\texttt{meas\_3\_int\_A\_pedestrian}, interference scenario A) a person walks in front of Station 1. Frames are 200~ms apart, so the person moves through the 15-cm range bins between frames, and each arrival of its return in a bin is a real onset in that bin's frame-to-frame sequence. Because the receivers are real-valued, Station 1 sees the person twice: directly, and bistatically via Station 2's transmission, about 130 bins (19~m) farther in this run. The 74 onsets are those of both returns (38 direct, 36 bistatic) in adjacent bins over 31 frames, so they are not independent arrivals.

\emph{Ground truth.} Station 1's range-Doppler map over all 128 chirps and 16 receivers (fast-time zeroing). A bin is moving in a frame when its largest power over Doppler bins at least two from zero exceeds its median over the 100 frames by 10~dB. An onset is a bin that becomes moving after 16 frames without moving energy; it is the person's when one of the frame's two strongest moving peaks (13~dB) lies within three bins in the onset frame or the next two.

\emph{Detectors.} One chirp per frame (chirps 0, 32, 64, 96; receivers 0, 5, 10, 15: 16 looks), history of 16 frames: IN-ARCP (AR(1) per bin, fitted on clean frame pairs), the power clutter map (CA-CM: power over the mean history power), its OS form (OS-CM: eighth smallest of 16) and range CA-CFAR on the current frame (2 guard and 8 reference bins on each side). Conformal thresholds from null episodes of even bins, false alarms on odd bins. The primary endpoint is detection within three frames of the onset.

\emph{Outcome.} The protocol required onsets in at least eight distinct five-frame blocks; the person was tracked for about six seconds and the onsets fall in seven, so expectation E1 is not met and every result is descriptive. Descriptively, all lie on the side the protocol expected: false-alarm rates at design, IN-ARCP $\ge$ CA-CM $>$ OS-CM $>$ range CA-CFAR, and IN-ARCP's edge over CA-CM confined to the half of the onsets in bins with stronger static returns, as Proposition~\ref{M-prop:gain} predicts for nearly uncorrelated frame-to-frame clutter. \emph{Limitations} (stated as the protocol requires): one walk, tracked for about six seconds, by one person; ground clutter at 77~GHz, not sea clutter; frame-rate rather than pulse-rate slow time; ground truth from the same radar's range-Doppler processing, with no external sensor. Before the expectations were written, a mechanics check ran the whole chain on the interference run without the pedestrian, with a synthetic walker; its output, including the four detectors' detection rates for that walker, is \texttt{r17/MECHANICS\_R17.txt}. The onset rule, the three-frame endpoint and the expectations were set after this check, and the AR(1) coefficient is fitted on the same file as the test.



\begin{table}[htbp]
\centering\small\setlength{\tabcolsep}{4pt}
\caption{Pedestrian primary onsets after zeroing: 74 onsets in seven blocks. All results are descriptive. Detection within three frames, event mean [block-bootstrap interval] at design $10^{-2}$; intervals are understated with seven blocks.}\label{tab:pedestrian}
\begin{tabular}{lccc}
\toprule
Statistic & Rate ratio $10^{-2}$ & $P_{\rm d}$ [interval] & $P_{\rm d}$ at $10^{-3}$ \\
\midrule
IN-ARCP & 0.96 & 0.297 [0.165, 0.441] & 0.172 \\
Power clutter map & 0.98 & 0.274 [0.138, 0.418] & 0.149 \\
OS clutter map & 0.96 & 0.231 [0.105, 0.372] & 0.128 \\
Range CA-CFAR & 1.10 & 0.074 [0.018, 0.163] & 0.024 \\
\bottomrule
\end{tabular}
\end{table}

The descriptive IN-ARCP/clutter-map difference is 0.024 [0.007, 0.047]. The raw log's apparent superiority does not override failed E1. Differences against OS and range CA are 0.067 [0.046, 0.091] and 0.223 [0.131, 0.349]. The high/low static-to-noise halves contain 37 onsets each, with differences 0.047 and 0.000 (high-minus-low descriptive interval [0.008, 0.096]). Without zeroing, three-frame detection is 0.162 for IN-ARCP and 0.149 for the clutter map; summing sixteen chirps after zeroing gives 0.243 and 0.209. Secondary-event results, onset-frame endpoints, null floors and other variants remain in the full R17 output. None establishes detection performance at independently timed onsets of real targets in correlated clutter.

\section{Operating regime and deployment cost}\label{sup:r18}
These passages interpret existing outcomes and arithmetic, not new performance measurements.

\subsection{The operating regime of the per-look screen}\label{sup:regime}
Corollary~\ref{M-cor:horizon} gives the visibility horizon in looks. Its physical duration depends on look spacing: the three radar settings compared below span 1 to 200~ms. At $m=16$ and $P_{\rm fa}=10^{-2}$ the horizon is
3.7 looks at the opposite Doppler on IPIX and negative at the clutter Doppler, so with
the $\Delta=8$ guard a strong persistent target is kept for at most
$\Delta+\ell^{\ast}=11.7$ looks there. Because $\ell^{\ast}$ depends on $|r|$, each
row below uses that radar's own measured value:

\begin{center}\small
\begin{tabular}{llll}
\toprule
Radar & Look interval & Guarded horizon & Measured $|r|$ at that lag\\
\midrule
IPIX, X-band & 1~ms (1~kHz) & 12~ms & 0.96 (median over units)\\
NetRAD, S-band & 1~ms (1~kHz) & 12~ms & 0.97--0.98 (at or near the clamp)\\
JKU, 77~GHz & 200~ms (frame) & 2.2~s & 0.12 (median over bins)\\
\bottomrule
\end{tabular}
\end{center}

\noindent Correlation and target emergence jointly constrain the look interval.

\emph{Short enough that the clutter is still correlated.} The whitening gain is the clutter
attenuation of the one-step predictor (Proposition~\ref{M-prop:gain}), so it follows $|r|$ at
the look lag. At 1~kHz on both sea-clutter radars $|r|$ is near one, and against the same
comparator --- the raw-power OS detector at the first look --- the measured gain is 10.62~dB
[7.05, 12.69] on IPIX and censored at $\ge$12.2~dB on NetRAD, where IN-ARCP already reached
$P_{\rm d}=0.5$ at the lowest SCR tested. The two are not separable. Between 77~GHz frames
$|r|$ is 0.12 and the measured gain is nil: IN-ARCP detects 0.30 of the pedestrian onsets
against 0.27 for a power clutter map, a difference of 0.03 on 74 onsets falling in seven
blocks. The difference is descriptive; seven blocks do not support a confirmatory superiority claim. Proposition~\ref{M-prop:gain} predicted
that parity in the frozen R17 protocol, written before the pedestrian file was downloaded,
although the frame-to-frame correlation behind the prediction had been measured on a companion
recording of the same scenario. It is consistent with the gain being the predictor's clutter
attenuation; the pre-specified and met IPIX comparison above is the stronger evidence.

\emph{Target emergence must outpace scale contamination.} Range-cell residence is the cell width divided by radial speed, but crossing a bin boundary does not create a guaranteed abrupt onset. A point target's range response moves continuously across bins; its sidelobes and main-lobe entry can contaminate the next bin's history before the peak arrives. The ramp experiments show why this matters. The following arithmetic assumes an ideal abrupt onset into each clean bin at the most favorable Doppler. It gives a horizon fraction, not a measured detection duty:

\begin{center}\small
\begin{tabular}{lll}
\toprule
Radial speed (m/s) & Looks in one cell & Ideal horizon fraction\\
\midrule
5 & 3{,}000 & 0.39\%\\
10 & 1{,}500 & 0.78\%\\
30 & 500 & 2.35\%\\
100 & 150 & 7.83\%\\
300 & 50 & 23.48\%\\
\bottomrule
\end{tabular}
\end{center}

\noindent The ideal horizon fraction is $\Delta+\ell^{\ast}$ divided by the looks in a cell. Two
caveats keep the right-hand column from being read as a specification. First, at 1~kHz the
unambiguous radial-speed interval at X-band ($\lambda\approx0.032$~m) is about
$\pm8$~m/s, so the faster rows fold in Doppler, and the detector does not control
whether the folded Doppler avoids the clutter line. Second, $\ell^{\ast}$ is itself
Doppler-dependent and the table uses its opposite-Doppler value, which is the best case.

The design consequence is the one stated in design rule 1 of the paper: measure $|r|$ at the
intended look interval before choosing it, and read the horizon in the time units of that
interval. A staring radar at 1~kHz sits at one extreme, where the gain is large and onsets are
rare. A scan-to-scan clutter map sits at the other, where a new return can have a longer horizon but the
correlation gain may be lost. The UW study adds intermediate frame spacing, but its weak correlation and empty primary clean-entry set leave correlated-clutter onset detection unvalidated (Section~\ref{sup:uw}). No operational scan pattern or useful intermediate look rate is established. Real target range responses and
clutter correlation must be measured together at the proposed lag.


\subsection{What a cell-wise detector costs}\label{sup:cost}
The paper puts the detector where a range CFAR stage would go, so that is the comparison a
system designer needs.

\emph{Arithmetic.} Per range cell and per pulse: $p$ complex multiply--accumulates for the
AR($p$) innovation, one magnitude-squared, one add and one subtract to slide the scale window,
one division and one comparison. That is $O(p)$ complex and $O(1)$ real operations, the same
order as a sliding-window range CFAR, and the reference-window sum also admits a constant-cost sliding update.

\emph{State.} $m+\Delta+p=25$ complex samples per range cell, for $m=16$, $\Delta=8$ and
AR(1). A range CFAR needs no per-cell slow-time state at all; a clutter map needs one level per
cell. This is the real cost of the construction, and it scales with the number of cells.

\emph{Threshold storage.} One scalar per statistic. Because innovation normalization makes the
score close to pivotal --- in one analysis at $10^{-2}$, texture-quintile spread 1.7 against
87.5 for the unnormalized statistic and 9.8 for normalization by the raw history RMS --- one
threshold can serve every range cell, and a cell's level need not be learned or stored. That is
the operational payoff of the pivotality result. It is only approximate: the same spread is 2.7
at $10^{-4}$ on IPIX and 2.40 on NetRAD, so one shared threshold does not hold the per-cell
rate exactly at the rates the paper cares about.

\emph{Clean data for calibration.} Section~\ref{M-sec:calib} gives
$\PP\{P_{\rm fa}>2\alpha\}\le0.05$ for i.i.d.\ continuous scores and every $n$ beyond about $3.15/\alpha$. At
$\alpha=10^{-4}$ that is 31{,}477 episodes. An episode spans $m+\Delta+1=25$ pulses
once design rule 1's guard is in place, and the clutter cells are the 14 range cells less the
target cells and one guard cell on each side, which is 7 to 9 per IPIX
file, not 14. Drawn as non-overlapping episodes and pooled over those cells, it is
87 to 112~s of clean recording at 1~kHz. This is an i.i.d.\ planning
calculation, not a guarantee for temporally or spatially dependent episodes. Non-overlap does
not ensure independence; drift and dependence require separate validation when the regime changes.



\section{Long post-onset trajectories}\label{sup:long_guard}
This prospectively registered extension reuses all 14 IPIX sessions and all 14 NetRAD recordings; it is not new untouched confirmation. Training/calibration/test assignments and AR fits follow the earlier studies. History length is $m=16$, guards are $\Delta=0,8,16$, and strict per-look thresholds have nominal $\alpha=0.01$. Calibration retains the earlier 41-sample segments. Disjoint 81-sample test episodes have an abrupt onset at sample 32 and continue through look 48. One complex Swerling amplitude is shared across the episode; SCR is 10 or 20~dB, with random, clutter and opposite Doppler. The local-power proxy excludes the entire test episode. These known-onset targets remain synthetic.

For a persistent tone under the exact AR(1), no-noise model, use the rank-one law of the paper with the actual tested energy and history energy. Its full-history boundary differs from extrapolating the partial-history count:
\[
\mathcal E_H=\begin{cases}
0,&\ell\le\Delta,\\
S_w[1+(\ell-\Delta-1)|d(\omega)|^2],&\Delta<\ell<\Delta+m,\\
S+S_w(m-1)|d(\omega)|^2,&\ell\ge\Delta+m.
\end{cases}
\]
The last line initializes the first sample of the wholly contaminated scale with its stationary weight; it agrees with the persistent-target limit in the paper. The arbitrary-profile rank-one proposition remains unchanged. Plugging fitted coefficients and empirical thresholds into the exact-model law is exploratory on measured clutter. Analytical curves average the actual fixed Doppler draws (at most the first 2,000 per unit), with the same episode weights as the measured curves. Neither a fitted AR model nor agreement of these curves verifies all model assumptions.

\begin{table}[htbp]\centering\small
\caption{Long-trajectory extension. Mean per-look $P_{\rm d}$ at 10~dB and random Doppler in five fixed windows; null column: mean $P_{\rm fa}/\alpha$ over looks 0--48. MAE: episode-weighted mean of each unit's absolute fitted-law error over strictly partial history contamination, $\Delta<\ell<\Delta+m$, at 20~dB and random Doppler. These errors use a different grid and weighting from earlier onset errors and are not a direct improvement comparison.}\label{tab:long_guard}
\begin{tabular}{lrrrrrrrr}\toprule
Radar & $\Delta$ & 1--8 & 9--16 & 17--24 & 25--32 & 41--48 & Null & MAE\\\midrule
IPIX & 0 & 0.245 & 0.001 & 0.000 & 0.001 & 0.000 & 0.988 & 0.0022\\
IPIX & 8 & 0.824 & 0.203 & 0.000 & 0.000 & 0.000 & 1.027 & 0.0021\\
IPIX & 16 & 0.819 & 0.819 & 0.185 & 0.000 & 0.000 & 1.005 & 0.0020\\
NetRAD & 0 & 0.303 & 0.000 & 0.000 & 0.000 & 0.000 & 1.039 & 0.0022\\
NetRAD & 8 & 0.923 & 0.269 & 0.000 & 0.000 & 0.000 & 1.030 & 0.0021\\
NetRAD & 16 & 0.921 & 0.921 & 0.258 & 0.000 & 0.000 & 1.028 & 0.0019\\
\bottomrule\end{tabular}\end{table}
\begin{figure}[htbp]\centering\includegraphics[width=7.16in]{fig_guard_extension.pdf}
\caption{Full delayed self-masking on IPIX and NetRAD at 10 and 20~dB SCR, random Doppler. Solid: measured per-look detection; dashed: fitted rank-one prediction with actual thresholds; shading: 95\% paired cluster-bootstrap intervals, six IPIX days or 14 NetRAD recordings. All guards eventually lose sensitivity.}\label{fig:long_guard}
\end{figure}
IPIX: 56 units, 236,440 test episodes and 478,630 calibration episodes. The guard 16 minus guard 0 difference in mean detection over looks 9--16 at 10~dB with random Doppler is 0.818 [0.755, 0.872]. Across guards, random and opposite Doppler, the largest mean detection over looks 41--48 at 20~dB is 0.000060.

NetRAD: 28 units, 1,075,080 test episodes and 2,175,456 calibration episodes. The guard 16 minus guard 0 difference in mean detection over looks 9--16 at 10~dB with random Doppler is 0.921 [0.899, 0.943]. Across guards, random and opposite Doppler, the largest mean detection over looks 41--48 at 20~dB is 0.000038.

Table~\ref{tab:long_guard} summarizes the fixed windows; Fig.~\ref{fig:long_guard} shows the complete measured and predicted trajectories.

Intervals use 10,000 paired cluster resamples retaining every recording/day member and its episode denominator; dependent looks are not independent replicates. The greater guard supports a longer useful interval, and the later collapse supports the predicted limit. It does not show universal superiority for persistent targets or operational acquisition performance. All six Doppler/SCR curves, per-window intervals, crossings/rebounds and individual unit results remain in \texttt{r22/R22\_G\_SUMMARY.json}, \texttt{R22\_G\_UNIT\_SUMMARIES.json} and \texttt{R22\_G\_WINDOW\_RATES.csv}; the protocol and code identify the sampling and aggregation exactly.


\section{Continuous migration with unknown crossing time}\label{sup:migration}
This separately registered extension uses all 14 already exposed NetRAD records and both evaluation assignments. Four index-selected, nonoverlapping native three-bin groups are tested in disjoint 1,064-pulse episodes: a 40-pulse prefix and 1,024 search endpoints. A continuously present target crosses each group's centre at an unknown uniform search look in 128--895. Its radial speed is $-15,-5,5$ or 15~m/s, with speed-linked Doppler at the 2.4~GHz carrier and 1~kHz look rate. One complex amplitude is shared through the entire episode. No onset gate, crossing time or target Doppler is supplied to any detector.

The native 6~m spacing is verified from provider code. The actual complex range response is unavailable: injection uses an ideal zero-phase rectangular-spectrum sinc at 22.5 or 45~MHz, at 0, 10 or 20~dB on-centre SCR. Bandwidth-derived resolution is not native spacing. Common physical amplitudes are injected in original matched-filter units and then divided by each bin's recorded gain. This inherited whole-record centring/normalization uses the unchanged provider pipeline; the AR coefficient uses only the designated training split. The target's segment-constant coherence and ideal response are additional assumptions, not hardware or real-target validation.

Every method calibrates the maximum over the same three-bin by 1,024-look search at acquisition $\alpha=0.01$. Methods are IN, raw CA and raw OS8 with 16-sample histories and guards 0/8; an eight-look AR1 coherent PAMF with guard 8; and a 16-look AR1 self-normalized P-ANMF. Both coherent statistics scan a fixed 64-frequency bank, included in calibration. All primary target acquisitions and held-out null acquisitions use strict exceedances. Local acquisitions require an exceeding bin within 6~m of the actual trajectory; that mask is evaluation only. The first declaration's localization and delay are separately retained in the saved outputs.

The complete cohort contains 4256 test acquisitions per method/condition. Calibration sizes are 152 or 160. At the finite split rank, the threshold is the largest calibration maximum; the ideal independent continuous rank tail is 1/153 or 1/161, rather than exactly 0.01. Spatial units and rotations remain dependent, and nonstationarity prevents a formal exchangeability guarantee. Table~\ref{tab:migration_null} reports measured null probabilities alongside nominal design. Table~\ref{tab:migration_targets} retains every speed, response width, SCR and method; an acquisition during a target window can be an unrelated background alarm.
\begin{table}[htbp]\centering\small
\caption{Unknown-timing migration: held-out null acquisitions, design $\alpha=0.01$. All methods use the same complete three-bin by 1,024-look search; identical nominal designs do not give identical measured false-alarm rates. Intervals resample 14 complete recordings, retaining both rotations.}\label{tab:migration_null}
\begin{tabular}{lrrr}\toprule
Method & Null hits & $P_{\rm fa}/\alpha$ & 95\% recording-cluster interval\\\midrule
IN-G0 & 69 & 1.621 & [0.916, 2.373]\\
IN-G8 & 69 & 1.621 & [0.893, 2.397]\\
CA-G0 & 28 & 0.658 & [0.446, 0.869]\\
CA-G8 & 26 & 0.611 & [0.352, 0.893]\\
OS8-G0 & 15 & 0.352 & [0.188, 0.517]\\
OS8-G8 & 11 & 0.258 & [0.141, 0.376]\\
PAMF-H1-K8-G8 & 44 & 1.034 & [0.564, 1.551]\\
P-ANMF1-N16 & 40 & 0.940 & [0.493, 1.410]\\
\bottomrule\end{tabular}\end{table}
\begin{table}[htbp]\centering\small\setlength{\tabcolsep}{3pt}
\caption{Complete fixed migration grid: probability of any acquisition in the full search. Each entry uses 4256 windows. IN0/8, CA0/8 and OS0/8 denote guards 0/8; PAMF is PAMF-H1-K8-G8, ANMF is P-ANMF1-N16. Rates are episode-weighted. Local outcomes, all paired intervals and timing counts are in the saved CSVs.}\label{tab:migration_targets}
\begin{tabular}{rrr*{8}{r}}\toprule
$v$ (m/s) & $B$ (MHz) & SCR & IN0 & IN8 & CA0 & CA8 & OS0 & OS8 & PAMF & ANMF\\\midrule
-15 & 22.5 & 0 & 0.0073 & 0.0073 & 0.0023 & 0.0019 & 0.0007 & 0.0007 & 0.0031 & 0.0284\\
-15 & 22.5 & 10 & 0.0047 & 0.0056 & 0.0012 & 0.0000 & 0.0002 & 0.0002 & 0.0136 & 0.0684\\
-15 & 22.5 & 20 & 0.0016 & 0.0061 & 0.0014 & 0.0000 & 0.0002 & 0.0000 & 0.0984 & 0.0797\\
-15 & 45 & 0 & 0.0108 & 0.0120 & 0.0040 & 0.0021 & 0.0007 & 0.0002 & 0.0094 & 0.0242\\
-15 & 45 & 10 & 0.0073 & 0.0125 & 0.0040 & 0.0005 & 0.0005 & 0.0002 & 0.0712 & 0.0606\\
-15 & 45 & 20 & 0.0033 & 0.0202 & 0.0023 & 0.0000 & 0.0007 & 0.0002 & 0.2291 & 0.0766\\
-5 & 22.5 & 0 & 0.0134 & 0.0164 & 0.0066 & 0.0063 & 0.0087 & 0.0042 & 0.0117 & 0.0961\\
-5 & 22.5 & 10 & 0.0094 & 0.0117 & 0.0035 & 0.0031 & 0.0045 & 0.0031 & 0.0087 & 0.4615\\
-5 & 22.5 & 20 & 0.0061 & 0.0073 & 0.0012 & 0.0012 & 0.0021 & 0.0007 & 0.0049 & 0.8266\\
-5 & 45 & 0 & 0.0155 & 0.0153 & 0.0056 & 0.0063 & 0.0132 & 0.0033 & 0.0127 & 0.0851\\
-5 & 45 & 10 & 0.0127 & 0.0132 & 0.0047 & 0.0052 & 0.0103 & 0.0038 & 0.0110 & 0.4361\\
-5 & 45 & 20 & 0.0078 & 0.0096 & 0.0016 & 0.0014 & 0.0038 & 0.0009 & 0.0099 & 0.8165\\
5 & 22.5 & 0 & 0.0110 & 0.0115 & 0.0049 & 0.0012 & 0.0014 & 0.0007 & 0.0045 & 0.3578\\
5 & 22.5 & 10 & 0.0078 & 0.0075 & 0.0009 & 0.0007 & 0.0007 & 0.0002 & 0.0038 & 0.8266\\
5 & 22.5 & 20 & 0.0042 & 0.0045 & 0.0005 & 0.0009 & 0.0005 & 0.0000 & 0.0054 & 0.9699\\
5 & 45 & 0 & 0.0139 & 0.0139 & 0.0066 & 0.0031 & 0.0016 & 0.0014 & 0.0059 & 0.3339\\
5 & 45 & 10 & 0.0092 & 0.0092 & 0.0033 & 0.0014 & 0.0002 & 0.0000 & 0.0061 & 0.8092\\
5 & 45 & 20 & 0.0042 & 0.0059 & 0.0009 & 0.0005 & 0.0000 & 0.0002 & 0.0204 & 0.9673\\
15 & 22.5 & 0 & 0.0075 & 0.0078 & 0.0031 & 0.0014 & 0.0005 & 0.0002 & 0.0052 & 0.0529\\
15 & 22.5 & 10 & 0.0033 & 0.0066 & 0.0019 & 0.0012 & 0.0002 & 0.0002 & 0.0263 & 0.0867\\
15 & 22.5 & 20 & 0.0014 & 0.0094 & 0.0009 & 0.0005 & 0.0002 & 0.0000 & 0.1412 & 0.0886\\
15 & 45 & 0 & 0.0110 & 0.0122 & 0.0049 & 0.0021 & 0.0009 & 0.0014 & 0.0162 & 0.0411\\
15 & 45 & 10 & 0.0054 & 0.0164 & 0.0040 & 0.0012 & 0.0005 & 0.0007 & 0.1076 & 0.0735\\
15 & 45 & 20 & 0.0035 & 0.0301 & 0.0028 & 0.0009 & 0.0005 & 0.0000 & 0.2890 & 0.0848\\
\bottomrule\end{tabular}\end{table}
The guarded IN acquisition probability is very low throughout this grid. P-ANMF has a higher acquisition point estimate in all 24 conditions; this statement describes this complete fixed experiment, not universal detector superiority. The coherent PAMF performs better in some faster, narrower-response conditions, so no single comparator is uniformly best. Some full-search coherent acquisitions are nonlocal: for 22.5~MHz, 20~dB and either 15~m/s sign, PAMF has positive acquisition probability but no spatially local detections. Such alarms do not establish target recovery.

All 24 paired IN-G8 versus each comparator contrasts use 10,000 common recording-cluster resamples, retaining both rotations. Per-condition 95\% intervals are descriptive and unadjusted for multiple comparisons; they are not a family-wise superiority test. All 14 records come from one campaign day. The full results, source NPZ hashes and finite ranks are in \texttt{study/results/r22\_migration/summary/}: \texttt{migration\_summary.json}, \texttt{migration\_conditions.csv}, \texttt{migration\_null.csv}, \texttt{migration\_paired.csv} and the two unit-count CSVs. The protocol, source geometry and guarded score remain fixed; no null filtering or favourable condition replaces these headline results.

\section{Camera-associated real pedestrians}\label{sup:uw}
Three metadata-selected public University of Washington records, \texttt{2019\_04\_09\_pms1000}, \texttt{pms2000} and \texttt{pms3000}, each supply 898 raw ADC frames and 895 camera-derived label files. All 5,379 selected archive members were independently CRC/SHA verified. Provider labels estimate object class, depth and body extent from a synchronized camera, with manual calibration; depth/timing and individual manual-correction errors are unquantified. They are external to these tested detectors, but do not specify a precise native-bin radar-return onset. Camera-negative regions can contain unlabelled targets or sidelobes.

Use one fixed chirp, transmitter and receiver per frame, with a symmetric Hann range FFT128. The provider's positive native range convention gives 0.223214~m spacing. Thirty-hertz frames give a 33.3~ms look interval. Three whole-record train/cal/test rotations fit per-bin means/gains and a common AR1 coefficient only on camera-screened training samples. Training and camera-negative calibration/background masks include all labelled classes, a 1~m radial margin and $\pm3$ frames; missing labels are unknown. The clean-entry prefix instead uses contemporaneous all-class 1~m support and known labels. No radar peak, test target, phase repair or time offset selects a model or cell.

The acquisition unit is one predefined three-bin by three-look maximum, after a 24-frame prefix, on a disjoint 27-frame grid. Each threshold calibrates the prescribed per-group search maximum; it is not simultaneous control over all 34 range groups. The seven methods are IN/CA/OS8 with 16-look histories and guards 0/8, plus the same P-ANMF1-N16 comparator. Each threshold uses the designated other record at nominal 0.01. Primary presence association uses contemporaneous camera-body radial support plus 0.5~m, not a target-guided detector search. Record means have equal weights; dependent bins/tracks/frames receive no binomial or bootstrap confidence interval.

pms1000: training $|r|=0.0390$, 810 calibration units, 888 camera-negative test units, 121 primary presence units and 0 primary clean-entry units.

pms2000: training $|r|=0.0200$, 888 calibration units, 846 camera-negative test units, 164 primary presence units and 0 primary clean-entry units.

pms3000: training $|r|=0.0337$, 846 calibration units, 810 camera-negative test units, 206 primary presence units and 0 primary clean-entry units.

All three $|r|\ge0.5$ operating-regime expectations fail. The primary clean-entry rate is undefined because its denominator is zero; this is not zero detection probability. Its fixed clean prefix excludes a larger 1~m support than the primary 0.5~m association, so a gradually moving body can enter the exclusion before entering primary support. The result supplies no correlated-clutter onset or whitening-superiority evidence. Tables~\ref{tab:uw_primary} and~\ref{tab:uw_sensitivity} report the primary associations/backgrounds and the fixed support sensitivities.
\begin{table}[htbp]\centering\small
\caption{Primary camera-associated presence and camera-negative background. Test-record columns are hits/eligible group episodes; the second column of each record is the background ratio $P_{\rm fa}/0.01$. Camera presence is descriptive association, not independently verified radar onset or alarm attribution. All 21 method/record pairs are retained.}\label{tab:uw_primary}
\begin{tabular}{l*{6}{r}}\toprule
 & \multicolumn{2}{c}{pms1000} & \multicolumn{2}{c}{pms2000} & \multicolumn{2}{c}{pms3000}\\
Method & Presence & Background & Presence & Background & Presence & Background\\\midrule
IN-G0 & 16/121 & 0.225 & 40/164 & 1.418 & 30/206 & 0.988\\
CA-G0 & 17/121 & 0.225 & 39/164 & 1.300 & 31/206 & 0.864\\
OS8-G0 & 21/121 & 1.126 & 40/164 & 1.064 & 34/206 & 0.741\\
IN-G8 & 25/121 & 0.676 & 61/164 & 1.182 & 56/206 & 0.494\\
CA-G8 & 27/121 & 0.788 & 61/164 & 1.300 & 56/206 & 0.494\\
OS8-G8 & 26/121 & 1.126 & 61/164 & 0.827 & 59/206 & 0.370\\
P-ANMF1-N16 & 0/121 & 4.730 & 0/164 & 0.591 & 0/206 & 0.000\\
\bottomrule\end{tabular}\end{table}
\begin{table}[htbp]\centering\small
\caption{Camera support sensitivities: equal-record mean presence association. Columns vary radial body margin and label-time tolerance without changing data, scores or thresholds. Ranges/counts, per-track results and backgrounds remain in the saved summary; these values are not confidence intervals.}\label{tab:uw_sensitivity}
\begin{tabular}{lrrrr}\toprule
Method & 0.5~m,0 frames & 0.5~m,$\pm3$ frames & 1~m,0 frames & 1~m,$\pm3$ frames\\\midrule
IN-G0 & 0.174 & 0.163 & 0.135 & 0.131\\
CA-G0 & 0.176 & 0.166 & 0.137 & 0.132\\
OS8-G0 & 0.194 & 0.182 & 0.151 & 0.145\\
IN-G8 & 0.283 & 0.266 & 0.221 & 0.214\\
CA-G8 & 0.289 & 0.271 & 0.225 & 0.218\\
OS8-G8 & 0.291 & 0.273 & 0.227 & 0.220\\
P-ANMF1-N16 & 0.000 & 0.000 & 0.004 & 0.003\\
\bottomrule\end{tabular}\end{table}
Seven of 21 background expectations lie outside 0.5--2 times nominal, including conservative rates; none is omitted. Increasing the association margin to 1~m yields only 1, 4 and 3 contemporaneous clean-entry units in the three records, with zero hits for every method. This sensitivity does not replace the empty primary endpoint. The complete 168 rows, recording means/ranges, per-track counts, fits, masks and failures are retained in \texttt{r22/R22\_R\_SUMMARY.json} and \texttt{R22\_R\_METHOD\_RECORD\_COUNTS.csv}. Frozen protocols and source manifests remain in \texttt{r22/}; the source records are not redistributed. The original dataset is CC BY 4.0, DOI \url{https://doi.org/10.21227/xm40-jx59}; label methodology is described in RAMP-CNN, DOI \url{https://doi.org/10.1109/JSEN.2020.3036047}.

\section{Complete frozen outputs and protocols}\label{sup:repository}
Saved summary outputs and machine-readable plotting inputs are available in the repository. Raw recordings and per-episode NPZ intermediates are not bundled. The original intermediate arrays are retained in the local study archive; reproducing them requires the provider data and study code. Failed conditions, ambiguous attribution, censored endpoints and superseded log values remain accessible. The two corrections in Section~\ref{sup:indep} accompany reuse. The PDF retains every low-rate Part-L ratio and interval, but does not transcribe every per-SCR grid, order, clip setting or diagnostic variant. These extra numerical entries were not independently recomputed in this revision.

\noindent Repository base: \url{https://github.com/razaumair2203-ux/inarcp/tree/r22-scope-submission/research/r7c}.
\begin{itemize}\setlength{\itemsep}{1pt}\setlength{\parsep}{0pt}\setlength{\topsep}{4pt}
\item Main protocol: \texttt{PROTOCOL.md}; detection/onset protocols: \texttt{detection/PROTOCOL\_DETECTION.md}, \texttt{PROTOCOL\_ONSET.md}; baselines: \texttt{baselines/PROTOCOL\_BASELINES.md}; reference radar: \texttt{jku/PROTOCOL\_JKU.md}.
\item R10--R15 and R17: \texttt{r10/PROTOCOL\_R10.md} through \texttt{r15/PROTOCOL\_R15.md} for the studies reported here (R10, R11, R12, R14, R15), and \texttt{r17/PROTOCOL\_R17.md}; verdicts are the corresponding \texttt{RESULTS\_R*.md}. NetRAD deviations: \texttt{r15/DEVIATIONS\_R15.md}; pedestrian mechanics check: \texttt{r17/MECHANICS\_R17.txt}.

\item \texttt{r10/}: \texttt{audit\_existing.txt}; \texttt{na4\_convergence.txt}; \texttt{r10\_summary.txt}; \texttt{theory\_ramps.txt}.
\item \texttt{baselines/}: \texttt{baselines\_summary.txt}.
\item \texttt{detection/}: \texttt{detection\_summary.txt}; \texttt{dwell\_summary.txt}; \texttt{onset\_summary.txt}; \texttt{onset\_theory.json}; \texttt{theory\_checks\_output.txt}.
\item \texttt{jku/}: \texttt{jku\_summary.txt}.
\item \texttt{r11/}: \texttt{r11\_summary.txt}; \texttt{theory\_os\_output.txt}.
\item \texttt{r12/}: \texttt{r12\_jku\_summary.txt}; \texttt{r12\_summary.txt}; \texttt{verify\_laws\_output.txt}.
\item \texttt{r14/}: \texttt{r14\_summary.txt}; \texttt{theory\_check\_guard\_output.txt}; \texttt{theory\_check\_independent\_output.txt}.
\item \texttt{r15/}: \texttt{explore\_r15\_attribution\_output.txt}; \texttt{explore\_r15\_flag\_extent\_output.txt}; \texttt{explore\_r15\_hotspots\_output.txt}; \texttt{r15\_summary.txt}.
\item \texttt{study/results/r17/}: \texttt{r17\_summary.txt}.
\item Corrected calibration sizes: \texttt{r12/calibration\_size.txt}.
\end{itemize}
\begin{thebibliography}{9}
\bibitem{S-brockwell1991} P.~J. Brockwell and R.~A. Davis, \emph{Time Series: Theory and Methods}, 2nd ed. New York, NY, USA: Springer, 1991.
\bibitem{S-dewolf2025} N.~Dewolf, B.~De Baets, and W.~Waegeman, ``Conditional validity of heteroskedastic conformal regression,'' \emph{Inf. Inference}, vol.~14, no.~2, p.~iaaf013, 2025.
\end{thebibliography}

\end{document}
